\chapter{Restriction, time reversal, and the classification of \texorpdfstring{$L$}{L}-symbols}
\label{ch:mpu_cocycles_l_symbol_classification}

This chapter continues the study of the anomaly and the $L$-symbols. The
anomaly three-cochain of the restriction of a group of matrix product
operators to a subgroup is the restriction of the anomaly
(Lemma~\ref{thm:mpug_family_restriction_omega}), and a time-reversal symmetry
makes the anomaly cohomologous to a real three-cochain
(Lemma~\ref{lem:mpug_time_reversal_tensor_real_anomaly}). After comparing
circle and nonzero complex coefficients in degree two
(Lemma~\ref{thm:mpug_circle_complex_units_h2}), the chapter studies block
independence of $L$-symbols and classifies the $L$-symbols of permuted blocks
through the stabilizer of one block: up to action-tensor gauges, those
compatible with a fixed three-cochain are in bijection with the degree-two
cohomology of the stabilizer (Lemma~\ref{thm:mpug_stabilizer_torsor}). It ends
with positive generalized cocycles, which are coboundaries
(Lemma~\ref{thm:mpug_positive_generalized_cocycle}), and with the sign that
decides whether a normalized three-cocycle on the group of order two has
nontrivial class (Lemma~\ref{thm:mpug_z2_nontrivial_scalar_class}).

\section{Restriction, time reversal, and block independence}

\subsection{Restriction to subgroups}

\begin{definition}[Restriction of a group of matrix product operators]
    \label{def:mpug_family_restriction}
    \lean{MPOTensor.GroupFamily.comap,
        MPOTensor.GroupFamily.FusionData.comap}
    \leanok
    \uses{def:mpug_anomaly_three_cochain}
    For a group homomorphism $f\colon H\to G$ and a family
    $\{O_g\}_{g\in G}$ as in Definition~\ref{def:mpug_anomaly_three_cochain},
    the restricted family is $\{O_{f(a)}\}_{a\in H}$. A choice of fusion
    tensors $(F^<_{g,h},F^>_{g,h})$ of the family restricts to the fusion
    tensors $(F^<_{f(a),f(b)},F^>_{f(a),f(b)})$ of the restricted family,
    where the bond space of $O_{f(a)f(b)}$ is identified with that of
    $O_{f(ab)}$ through the equality $f(a)f(b)=f(ab)$.
\end{definition}

\begin{lemma}[Anomaly three-cochains of restrictions]
    \label{thm:mpug_family_restriction_omega}
    \lean{MPOTensor.GroupFamily.IsNormalRepresentation.comap,
        MPOTensor.GroupFamily.castMat_comap,
        MPOTensor.GroupFamily.ext_of_heq,
        MPOTensor.GroupFamily.FusionData.comap_leftV,
        MPOTensor.GroupFamily.FusionData.comap_rightV,
        MPOTensor.GroupFamily.FusionData.omega_comap,
        MPOTensor.GroupFamily.FusionData.cyclicInvariant_omega_map,
        MPOTensor.GroupFamily.FusionData.cyclicInvariant_omega_map_of_comap_eq}
    \leanok
    \uses{def:mpug_family_restriction, def:mpug_cochain_restriction,
        def:mpug_cyclic_invariant}
    The restriction of a family satisfying the hypotheses of
    Definition~\ref{def:mpug_anomaly_three_cochain} along $f\colon H\to G$
    satisfies them for $H$, and the anomaly three-cochain of the restricted
    fusion tensors is $f^*\omega$. Consequently, if $a\in H$ and $a^n=e$,
    then for every choice of fusion tensors of the restricted family, with
    anomaly three-cochain $\omega_H$,
    \begin{align}
        I_{f(a),n}(\omega) & =I_{a,n}(\omega_H).
    \end{align}
\end{lemma}

\begin{proof}\leanok
    \uses{def:mpug_family_restriction, thm:mpug_anomaly_three_cochain_spec,
        thm:mpug_anomaly_three_cocycle_gauge, thm:mpug_cyclic_invariant_gauge}
    The operator law of the restricted family is that of the family at the
    elements $f(a)$, $f(b)$ and $f(ab)=f(a)f(b)$. The two fusion trees of the
    restricted fusion tensors at $(a,b,c)$ are those of the family at
    $(f(a),f(b),f(c))$, followed by one identification of the bond space of
    $O_{f(a)f(b)f(c)}$ with that of $O_{f(abc)}$. Left multiplication by this
    permutation matrix preserves the dressed identity, so
    Lemma~\ref{thm:mpug_anomaly_three_cochain_spec} gives
    $\omega_H(a,b,c)=\omega(f(a),f(b),f(c))$ for the restricted fusion
    tensors. For any other choice of fusion tensors of the restricted family,
    Lemma~\ref{thm:mpug_anomaly_three_cocycle_gauge} makes the two
    three-cochains cohomologous, and
    Lemma~\ref{thm:mpug_cyclic_invariant_gauge} makes their cyclic
    invariants at $a$ equal. Finally $I_{a,n}(f^*\omega)=I_{f(a),n}(\omega)$
    because $f(a^k)=f(a)^k$.
\end{proof}

\begin{lemma}[Values of the anomaly three-cochain from finitely many identities]%
    \label{thm:mpug_omega_integer_certificate}
    \lean{MPSTensor.isDressedProportional_complexOfInt,
        MPSTensor.isDressedProportional_complexOfInt_of_mem,
        MPOTensor.GroupFamily.FusionData.omega_eq_one_of_leftV_eq_rightV,
        MPOTensor.GroupFamily.ofBondOne,
        MPOTensor.GroupFamily.FusionData.ofBondOne,
        MPOTensor.GroupFamily.FusionData.omega_ofBondOne}
    \leanok
    \uses{def:mpug_anomaly_three_cochain}
    Let $T$ be a tensor with integer letters, $A$ a family of integer
    matrices, $L$ and $R$ integer matrices, and $z\in\{1,-1\}$. Suppose
    that $LT^aT^b=A^aLT^b$ for all letters $a,b$ and that $LT^a=zRT^a$ for
    every letter $a$. Then $LT^{\mathbf w}=zRT^{\mathbf w}$ for every nonempty
    word $\mathbf w$. It suffices to check these identities on a list of pairs
    $(T^a,A^a)$ that contains the pair of every letter. In particular, for the
    anomaly three-cochain of Definition~\ref{def:mpug_anomaly_three_cochain}:
    if the two fusion trees of a triple coincide, then $\omega$ is one at
    that triple; and a family of bond-one tensors whose stacked products
    $O_gO_h$ equal $O_{gh}$ has, with all fusion tensors equal to one, the
    constant three-cochain one.
\end{lemma}

\begin{proof}\leanok
    \uses{def:mpug_anomaly_three_cochain, thm:mpug_anomaly_three_cochain_spec}
    Since $z^2=1$, the identity $LT^a=zRT^a$ gives $RT^a=zLT^a$, and hence
    $RT^aT^b=zA^aLT^b=A^aRT^b$. Induction on the word gives the dressed
    identity, and the entrywise image in $\C$ preserves all products. When
    the two fusion trees coincide, the scalar $1$ satisfies the dressed
    identity, so Lemma~\ref{thm:mpug_anomaly_three_cochain_spec} gives
    $\omega=1$. For bond dimension one, both fusion trees are products of
    one-by-one identity matrices and identifications of one-dimensional
    spaces, so they coincide.
\end{proof}

\subsection{Time reversal}

\begin{definition}[Unimodular square root and modulus of a phase]
    \label{def:mpug_phase_root}
    \lean{TNLean.Algebra.phaseRoot,
        TNLean.Algebra.modulusUnit}
    \leanok
    For $u\in\C^\times$, let $s_u=e^{i\arg(u)/2}$ be the unimodular square root
    of the phase of $u$, with $|u|$ viewed as an element of $\C^\times$.
\end{definition}

\begin{lemma}[Unimodular square roots of phases]
    \label{thm:mpug_phase_root}
    \lean{TNLean.Algebra.star_phaseRoot,
        TNLean.Algebra.phaseRoot_mul_phaseRoot_mul_modulusUnit,
        TNLean.Algebra.star_eq_self_of_star_eq_ofReal_mul}
    \leanok
    \uses{def:mpug_phase_root}
    For $u\in\C^\times$, $\overline{s_u}=s_u^{-1}$ and $u=s_u^2\,|u|$. If
    $w\in\C^\times$ and $\overline w=p\,w$ for a positive real number $p$,
    then $p=1$ and $w$ is real.
\end{lemma}

\begin{proof}\leanok
    The first two claims are $\overline{e^{i\theta}}=e^{-i\theta}$ and the
    polar form $u=|u|e^{i\arg u}$. For the third, taking absolute values in
    $\overline w=p\,w$ gives $p=1$.
\end{proof}

\begin{lemma}[Time reversal makes the anomaly real]
    \label{thm:mpug_time_reversal_real_anomaly}
    \lean{TNLean.Algebra.ScalarThreeCochain.coboundary_mul,
        TNLean.Algebra.ScalarThreeCochain.star_fusionGauge,
        TNLean.Algebra.ScalarThreeCochain.coboundary_modulusUnit_pos,
        TNLean.Algebra.ScalarThreeCochain.star_fusionGauge_phaseRoot,
        TNLean.Algebra.ScalarThreeCochain.exists_cohomologousTo_star_eq_self}
    \leanok
    \uses{def:mpug_fusion_gauge, def:mpug_cohomologous_to,
        def:mpug_phase_root}
    Let $\omega$ be a scalar $3$-cochain and $\beta$ a scalar $2$-cochain with
    $\beta\mathbin{\triangleright}\omega=\overline\omega$. Then
    $s\mathbin{\triangleright}\omega$ is real-valued, where
    $s(g,h)=s_{\beta(g,h)}$ is the unimodular square root of the phase of
    $\beta(g,h)$ from Definition~\ref{def:mpug_phase_root}. Consequently, a
    scalar $3$-cochain cohomologous to its complex conjugate is cohomologous
    to a real-valued one. This is the conclusion
    of~\cite[Section~4.3]{GarreRubio2023MPOSym} that, when time reversal
    commutes with a group matrix product operator symmetry, the three-cocycle
    can be gauge transformed to a real number; no finiteness of $G$ is
    needed.
    % Source anchor: arXiv:2203.12563, lines 873--891.
\end{lemma}

\begin{proof}\leanok
    \uses{def:mpug_fusion_gauge, thm:mpug_phase_root}
    The coboundary is multiplicative in $\beta$ and commutes with complex
    conjugation. Writing $\beta=s^2N$ with $N(g,h)=|\beta(g,h)|$ and using
    $\overline s=s^{-1}$ gives
    \begin{align}
        \overline{s\mathbin{\triangleright}\omega}
         & =s^{-1}\mathbin{\triangleright}\overline\omega
        =s^{-1}\mathbin{\triangleright}(s^2N\mathbin{\triangleright}\omega)
        =N\mathbin{\triangleright}(s\mathbin{\triangleright}\omega).
    \end{align}
    The values of $dN$ are positive real numbers, so
    Lemma~\ref{thm:mpug_phase_root} gives $dN=1$ and
    $\overline{s\mathbin{\triangleright}\omega}
        =s\mathbin{\triangleright}\omega$. If
    $\omega=\beta\mathbin{\triangleright}\overline\omega$, conjugating gives
    $\overline\beta\mathbin{\triangleright}\omega=\overline\omega$, and the
    first part applies.
\end{proof}

\begin{lemma}[Time reversal makes the anomaly and the L-symbols jointly real]
    \label{thm:mpug_time_reversal_real_l_symbols}
    \lean{TNLean.Algebra.LSymbol.star_gauge,
        TNLean.Algebra.LSymbol.gauge_modulusUnit_pos,
        TNLean.Algebra.LSymbol.star_gauge_phaseRoot,
        TNLean.Algebra.LSymbol.exists_gauge_star_eq_self}
    \leanok
    \uses{def:mpug_fusion_gauge, def:mpug_l_gauge, def:mpug_phase_root}
    Let $\omega$ be a scalar $3$-cochain, $L$ an $L$-symbol, $\beta$ a
    scalar $2$-cochain and $\gamma$ an action-tensor scalar gauge with
    $\beta\mathbin{\triangleright}\omega=\overline\omega$ and
    $(\beta,\gamma)\mathbin{\triangleright}L=\overline L$. Let $s$ and $t$ be
    the unimodular square roots of the phases of $\beta$ and $\gamma$. Then
    $s\mathbin{\triangleright}\omega$ and
    $(s,t)\mathbin{\triangleright}L$ are both real-valued. This result is not stated
    in~\cite[Section~4.3]{GarreRubio2023MPOSym}.
    The source derives $\beta\mathbin{\triangleright}\omega=\overline\omega$
    and, with the fusion-tensor phase omitted,
    $(1,\gamma)\mathbin{\triangleright}L=\overline L$; the hypothesis here,
    with the same $\beta$ in both relations, is the gauge-consistent pair
    and is assumed rather than derived from time reversal.
    % Source anchor: arXiv:2203.12563, lines 873--891 and 935--957.
\end{lemma}

\begin{proof}\leanok
    \uses{def:mpug_l_gauge, thm:mpug_l_gauge_composition,
        thm:mpug_time_reversal_real_anomaly, thm:mpug_phase_root}
    The claim for $\omega$ is Lemma~\ref{thm:mpug_time_reversal_real_anomaly}.
    For $L$, writing $\beta=s^2N$ and $\gamma=t^2M$ with
    $N_{g,h}=|\beta_{g,h}|$ and $M^x_g=|\gamma^x_g|$ gives, as in that proof,
    $\overline{(s,t)\mathbin{\triangleright}L}
        =(N,M)\mathbin{\triangleright}((s,t)\mathbin{\triangleright}L)$, whose
    prefactor
    $M^x_{gh}N_{g,h}/(M^{h x}_gM^x_h)$ is a positive real number, hence
    equal to one by Lemma~\ref{thm:mpug_phase_root}.
\end{proof}

\begin{lemma}[The L-symbol relation imposed by time reversal]
    \label{thm:mpug_time_reversal_l_symbol_relation}
    \lean{TNLean.Algebra.LSymbol.star_apply_eq_gauge_one_apply,
        TNLean.Algebra.LSymbol.phase_sq_eq_of_gauge_eq_star}
    \leanok
    \uses{def:mpug_l_gauge}
    Let $T\neq0$ and $P$ be vectors in a complex vector space and let
    $L^x_{g,h}$ and $\gamma$ be as in Definition~\ref{def:mpug_l_gauge}.
    If $P=\overline{L^x_{g,h}}\,T$ and
    $\gamma^x_h\gamma^{h x}_gP=L^x_{g,h}\gamma^x_{gh}T$, then
    \begin{align}
        \frac{\overline{L^x_{g,h}}}{L^x_{g,h}}
         & =\frac{\gamma^x_{gh}}{\gamma^x_h\gamma^{h x}_g}.
        \label{eq:mpug_time_reversal_l_symbol}
    \end{align}
    More generally, if $(\beta,\gamma)\mathbin{\triangleright}L=\overline L$,
    then for all $x$, $g$ and $h$,
    \begin{align}
        \left(\frac{L^x_{g,h}}{|L^x_{g,h}|}\right)^2
         & =\frac{\gamma^{h x}_g\gamma^x_h}{\gamma^x_{gh}\beta_{g,h}}.
    \end{align}
    In~\cite[Section~4.3]{GarreRubio2023MPOSym}, $T$ and $P$ are the
    complex conjugates of the two composite reductions compared by the
    $L$-symbol, the first hypothesis is the conjugate of the defining
    relation of $L$, and the second is the display derived from time
    reversal. The source prints the consequence as
    $|L^x_{g,h}|^2=\gamma^x_{gh}/(\gamma^x_h\gamma^{h x}_g)$; that form fails
    for an on-site $\mathbb Z_2\times\mathbb Z_2$ symmetry, and
    (\ref{eq:mpug_time_reversal_l_symbol}) is the local correction recorded
    in~\cite{gap:glm23_time_reversal_l_symbol}. The second display is the
    corrected sense in which the $L$-symbols squared are gauge
    transformations: the square of the phase of $L$ is, while its modulus is
    unconstrained.
    % Source anchor: arXiv:2203.12563, lines 935--957 (TRSF1group at
    %   line 954).
    % Local fix: docs/paper-gaps/glm23_time_reversal_l_symbol.tex.
\end{lemma}

\begin{proof}\leanok
    \uses{def:mpug_l_gauge}
    Substituting the first hypothesis into the second gives
    $\gamma^x_h\gamma^{h x}_g\overline{L^x_{g,h}}\,T
        =L^x_{g,h}\gamma^x_{gh}T$, and $T\neq0$ gives
    (\ref{eq:mpug_time_reversal_l_symbol}). For the second claim,
    $\overline L=cL$ with
    $c=\gamma^x_{gh}\beta_{g,h}/(\gamma^{h x}_g\gamma^x_h)$, so
    $|L|^2=\overline LL=cL^2$ and $(L/|L|)^2=c^{-1}$.
\end{proof}

\begin{definition}[The time-reversed tensor]
    \label{def:mpug_time_reversed_tensor}
    \lean{MPOTensor.timeReversal}
    \leanok
    Let $v$ be a $d\times d$ matrix and $T$ a matrix product operator tensor.
    The time-reversed tensor $\mathcal T T$ has letters
    $(\mathcal T T)^{ij}=\sum_{k,l}v_{ik}\,\overline{T^{kl}}\,\overline{v_{jl}}$,
    where the bar is entrywise complex conjugation of the virtual matrices.
    For the antiunitary time reversal $\mathcal T=v^{\otimes N}K$ of
    \cite[Section~4.3]{GarreRubio2023MPOSym}, with $K$ complex conjugation,
    it generates $\mathcal T O\mathcal T^{-1}$.
    % Source anchor: arXiv:2203.12563, line 867.
\end{definition}

\begin{lemma}[Operators and products of time-reversed tensors]
    \label{lem:mpug_time_reversed_tensor_properties}
    \lean{MPOTensor.mpo_map_starRingEnd, MPOTensor.mpo_timeReversal,
        MPOTensor.toMPSTensor_timeReversal, MPOTensor.mulTensor_timeReversal}
    \leanok
    \uses{def:mpug_time_reversed_tensor, def:mpdo_operator_product}
    For every $N$, the periodic operator of $\mathcal T T$ is
    $O_N(\mathcal T T)=v^{\otimes N}\,\overline{O_N(T)}\,(v^{\otimes N})^\dagger$.
    On the doubled physical index, each letter of $\mathcal T T$ is the
    combination $\sum_{k,l}v_{ik}\overline{v_{jl}}\,\overline{T^{kl}}$ of the
    conjugated letters of $T$. If $v^\dagger v=1$, then time reversal
    commutes with stacking: the time-reversed tensor of the stacked product
    of two tensors is the stacked product of their time-reversed tensors.
\end{lemma}

\begin{proof}
    \leanok
    The periodic operator is the trace of a product of letters, and complex
    conjugation commutes with products and traces. The physical factors
    $v$ and $\overline v$ act on the ket and bra legs and produce
    $v^{\otimes N}$ and $(\overline v^{\otimes N})^T=(v^{\otimes N})^\dagger$.
    For stacking, the contracted middle leg carries
    $\sum_k\overline{v_{kb}}v_{kc}=(v^\dagger v)_{bc}=\delta_{bc}$.
\end{proof}

\begin{lemma}[Reductions under recombination, conjugation and gauges]
    \label{lem:mpug_reduction_recombination}
    \lean{MPSTensor.IsReduction.recombine, MPSTensor.IsReduction.map_starRingEnd,
        MPSTensor.IsReduction.target_gauge,
        MPSTensor.IsDressedProportional.recombine,
        MPSTensor.IsDressedProportional.map_starRingEnd,
        MPSTensor.IsReduction.timeReversal,
        MPSTensor.IsDressedProportional.timeReversal}
    \leanok
    \uses{def:mps_rectangular_reduction, thm:mpug_dressed_proportionality,
        lem:mpug_time_reversed_tensor_properties}
    Let $(V,W)$ reduce $B$ onto $A$. Replacing the letters of both tensors by
    the same combinations $\sum_ac_{ia}B^a$ and $\sum_ac_{ia}A^a$ gives a
    reduction with the same $V,W$; conjugating the tensors and both
    boundaries gives a reduction; and if $A'^iQ=QA^i$ with $PQ=1$, a
    reduction $(V,W)$ of $B$ onto $A'$ gives the reduction $(PV,WQ)$ onto $A$.
    If $XB^{\mathbf w}=z\,YB^{\mathbf w}$ for long words, the same holds for
    recombined letters, and conjugation gives
    $\overline X\,\overline B^{\mathbf w}=\overline z\,\overline Y\,\overline B^{\mathbf w}$.
    In particular, reductions and dressed proportionalities of tensors pass
    to their time-reversed tensors with conjugated boundaries and scalars.
\end{lemma}

\begin{proof}
    \leanok
    A word in the recombined letters is the combination, with the same
    coefficients for both tensors, of words of the same length in the
    original letters; induction on the word with an arbitrary prefix gives
    both recombination claims. Conjugation commutes with products. For the
    gauge, $A'^{\mathbf w}Q=QA^{\mathbf w}$ gives
    $PVB^{\mathbf w}WQ=PA'^{\mathbf w}Q=A^{\mathbf w}$.
\end{proof}

\begin{lemma}[Fusion tensors along exact bond gauges]
    \label{lem:mpug_fusion_tensor_gauge_transport}
    \lean{MPOTensor.bondKron, MPOTensor.bondKron_eq_idKron_mul_kronId,
        MPOTensor.bondKron_mul_bondKron, MPOTensor.bondKron_one, MPOTensor.bondKron_mul_kronId,
        MPOTensor.bondKron_mul_idKron, MPOTensor.bondKron_bondKron_mul_assocInv,
        MPOTensor.mulTensor_mul_bondKron_of_intertwine,
        MPOTensor.GroupFamily.withTensor, MPOTensor.GroupFamily.castMat_mul_eq,
        MPOTensor.GroupFamily.BondGauge, MPOTensor.GroupFamily.BondGauge.mul_P,
        MPOTensor.GroupFamily.FusionData.gaugeTransport,
        MPOTensor.GroupFamily.FusionData.gaugeTransport_leftV,
        MPOTensor.GroupFamily.FusionData.gaugeTransport_rightV,
        MPOTensor.GroupFamily.FusionData.isAssociator_gaugeTransport}
    \leanok
    \uses{def:mpug_anomaly_three_cochain, lem:mpug_reduction_recombination,
        thm:mpug_reduction_dagger}
    Let $g\mapsto T_g$ be tensors with the bond dimensions of a group family
    $g\mapsto O_g$, and let $P_g,Q_g$ be mutually inverse matrices with
    $T_g^{il}Q_g=Q_gO_g^{il}$. If $(V_{g,h},W_{g,h})$ are fusion tensors of
    the family $T$, then
    \begin{align}
        V'_{g,h} & =P_{gh}V_{g,h}(Q_g\otimes Q_h), &
        W'_{g,h} & =(P_g\otimes P_h)W_{g,h}Q_{gh}
        \notag
    \end{align}
    are fusion tensors of the family $O$, and the two fusion trees of the
    triple product compare by the same scalar for both choices.
\end{lemma}

\begin{proof}
    \leanok
    The gauges $P_g\otimes P_h$ intertwine the stacked letters, so
    Lemma~\ref{lem:mpug_reduction_recombination} gives the reductions. The
    tensor products of the gauges commute with identity strands and with the
    bond associator, which gives
    $V'^L=P_{ghk}V^L(Q_g\otimes Q_h\otimes Q_k)$ and the same for the second
    tree. The gauge $Q_g\otimes Q_h\otimes Q_k$ intertwines the letters of
    the two triple products, so the dressed proportionality of the trees
    is transported unchanged.
\end{proof}

\begin{lemma}[Conjugated fusion tensors of the time-reversed family]
    \label{lem:mpug_conjugated_fusion_tensors}
    \lean{MPOTensor.kronId_map_starRingEnd, MPOTensor.idKron_map_starRingEnd,
        MPOTensor.mulTensorAssocInvMatrix_map_starRingEnd,
        MPOTensor.GroupFamily.castMat_map_starRingEnd,
        MPOTensor.GroupFamily.FusionData.timeReversal,
        MPOTensor.GroupFamily.FusionData.timeReversal_leftV,
        MPOTensor.GroupFamily.FusionData.timeReversal_rightV,
        MPOTensor.GroupFamily.FusionData.isAssociator_timeReversal}
    \leanok
    \uses{def:mpug_anomaly_three_cochain,
        lem:mpug_time_reversed_tensor_properties,
        lem:mpug_reduction_recombination}
    Let $v^\dagger v=1$ and let $(V_{g,h},W_{g,h})$ be fusion tensors of a
    group family $O$. Then $(\overline{V_{g,h}},\overline{W_{g,h}})$ are
    fusion tensors of the time-reversed family $g\mapsto\mathcal TO_g$, and
    if the fusion trees of $O$ compare by $z$, those of the conjugated
    fusion tensors compare by $\overline z$.
\end{lemma}

\begin{proof}
    \leanok
    Time reversal commutes with stacking, by
    Lemma~\ref{lem:mpug_time_reversed_tensor_properties}, so the stacked
    and triple products of the time-reversed family are the time-reversed
    tensors of those of $O$. The identity strands, bond identifications and
    the bond associator are real, so the fusion trees of the conjugated
    tensors are the conjugates of the original trees, and
    Lemma~\ref{lem:mpug_reduction_recombination} applies.
\end{proof}

\begin{lemma}[The gauges of a time-reversal symmetry]
    \label{lem:mpug_time_reversal_bond_gauge}
    \lean{MPOTensor.GroupFamily.CommutesWithTimeReversal,
        MPOTensor.GroupFamily.nonempty_bondGauge_timeReversal}
    \leanok
    \uses{def:mpug_anomaly_three_cochain, cor:asym_single_block,
        lem:mpug_time_reversed_tensor_properties}
    Let $g\mapsto O_g$ be a group of matrix product operators with normal
    tensors, let $v$ be unitary, and suppose that
    $v^{\otimes N}\overline{O_g}=O_gv^{\otimes N}$ for every $g$ and every
    $N\geq1$, that is, $\mathcal TO_g=O_g\mathcal T$. Then for every $g$
    there are mutually inverse matrices $P_g,Q_g$ with
    $(\mathcal TT_g)^{il}Q_g=Q_gT_g^{il}$, where $T_g$ is the tensor of
    $O_g$. These are the matrices $q_g$ of
    \cite[Section~4.3]{GarreRubio2023MPOSym}.
    % Source anchor: arXiv:2203.12563, lines 869--870.
\end{lemma}

\begin{proof}
    \leanok
    By Lemma~\ref{lem:mpug_time_reversed_tensor_properties} and
    unitarity, $\mathcal TT_g$ and $T_g$ generate the same periodic
    operators. Corollary~\ref{cor:asym_single_block} gives a reduction
    $(P_g,Q_g)$ of $\mathcal TT_g$ onto the normal tensor $T_g$. The bond
    dimensions agree, so $P_gQ_g=1$ forces $Q_gP_g=1$, and the one-letter
    words give $(\mathcal TT_g)^{il}=Q_gT_g^{il}P_g$.
\end{proof}

\begin{lemma}[Time reversal makes the anomaly of a group of operators real]
    \label{lem:mpug_time_reversal_tensor_real_anomaly}
    \lean{MPOTensor.GroupFamily.FusionData.omega_cohomologousTo_star_of_timeReversal,
        MPOTensor.GroupFamily.FusionData.exists_real_cohomologousTo_omega_of_timeReversal}
    \leanok
    \uses{lem:mpug_time_reversal_bond_gauge,
        lem:mpug_conjugated_fusion_tensors,
        lem:mpug_fusion_tensor_gauge_transport,
        thm:mpug_anomaly_three_cocycle_gauge,
        thm:mpug_time_reversal_real_anomaly}
    Under the hypotheses of Lemma~\ref{lem:mpug_time_reversal_bond_gauge},
    the anomaly three-cochain $\omega$ of any choice of fusion tensors is
    cohomologous to its complex conjugate, and hence to a real-valued
    three-cochain. This is the conclusion
    of~\cite[Section~4.3]{GarreRubio2023MPOSym} that
    $\omega=\overline\omega\,\beta_{g,h}\beta_{gh,k}/(\beta_{h,k}\beta_{g,hk})$
    and that $\omega$ can be gauge transformed to a real number.
    % Source anchor: arXiv:2203.12563, lines 869--891.
\end{lemma}

\begin{proof}
    \leanok
    Conjugate the fusion tensors, by
    Lemma~\ref{lem:mpug_conjugated_fusion_tensors}, and transport them
    along the gauges of Lemma~\ref{lem:mpug_time_reversal_bond_gauge}, by
    Lemma~\ref{lem:mpug_fusion_tensor_gauge_transport}. This gives fusion
    tensors of the original family whose trees compare by
    $\overline{\omega}$, so their anomaly three-cochain is $\overline\omega$.
    By Lemma~\ref{thm:mpug_anomaly_three_cocycle_gauge}, the two choices of
    fusion tensors have cohomologous three-cochains; the scalars $\beta_{g,h}$
    of the source are their relative scalars. The second claim is
    Lemma~\ref{thm:mpug_time_reversal_real_anomaly}.
\end{proof}

\subsection{Pullback and extension of degree-two classes}

\begin{definition}[Pullback of scalar cohomology]
    \label{def:mpug_h2_pullback}
    \lean{TNLean.Algebra.ScalarCocycle.comap,
        TNLean.Algebra.ScalarCocycle.comap_apply,
        TNLean.Algebra.ScalarCocycle.mapCocycles₂_cocyclesOfIsMulCocycle₂,
        TNLean.Algebra.ScalarCocycle.map_H2π}
    \leanok
    \uses{def:h2_cx}
    Let $f:K\to G$ be a group homomorphism. For a scalar $2$-cochain
    $\omega:G\times G\to\C^\times$, define
    $(f^*\omega)(a,b)=\omega(f(a),f(b))$.
    The restriction map
    $f^*:H^2(G,\C^\times)\to H^2(K,\C^\times)$ induced by $f$, with the
    identity of $\C^\times$ as coefficient map, satisfies
    $f^*[\omega]=[f^*\omega]$ for every $2$-cocycle $\omega$.
\end{definition}

\begin{lemma}[Pullback preserves cocycles and cohomology]
    \label{thm:mpug_h2_pullback_preserves}
    \lean{TNLean.Algebra.ScalarCocycle.isMulCocycle₂_comap,
        TNLean.Algebra.isMulCoboundary₂_div_comap}
    \leanok
    \uses{def:mpug_h2_pullback, def:cohomologous_to, def:is_cocycle}
    Let $f:K\to G$ be a group homomorphism. If $\omega$ is a scalar
    $2$-cocycle on $G$, then $f^*\omega$ is a scalar $2$-cocycle on $K$.
    If $\omega_1$ and $\omega_2$ are cohomologous scalar $2$-cochains on $G$,
    then $f^*\omega_1$ and $f^*\omega_2$ are cohomologous on $K$.
\end{lemma}

\begin{proof}\leanok
    \uses{def:mpug_h2_pullback, def:cohomologous_to, def:is_cocycle}
    For $a,b,c\in K$, substituting $f(a),f(b),f(c)$ into the cocycle equation
    for $\omega$ and using $f(ab)=f(a)f(b)$ and $f(bc)=f(b)f(c)$ shows that
    $f^*\omega$ is a $2$-cocycle.
    If $\omega_1(g,h)=\varphi(g)\varphi(h)\varphi(gh)^{-1}\omega_2(g,h)$,
    then with $\eta=\varphi\circ f$, substituting $g=f(a)$ and $h=f(b)$ gives
    $(f^*\omega_1)(a,b)=\eta(a)\eta(b)\eta(ab)^{-1}(f^*\omega_2)(a,b)$,
    so $\eta$ gives the required cohomology relation.
\end{proof}

\begin{lemma}[Restriction is the inhomogeneous restriction map]
    \label{lem:mpug_restriction_inhomogeneous}
    \lean{TNLean.Algebra.scalarRepresentationRes,
        TNLean.Algebra.ScalarCocycle.toInhomogeneousCochain_comap,
        TNLean.Algebra.ScalarThreeCochain.toInhomogeneousCochain_comap}
    \leanok
    \uses{def:mpug_h2_pullback, def:mpug_cochain_restriction,
        def:mpug_h3_group_cohomology}
    Let $f:K\to G$ be a group homomorphism. Under the identification of scalar
    $2$- and $3$-cochains with inhomogeneous cochains valued in the trivial
    representation $\C^\times_{\mathrm{triv}}$, the pullback $f^*$ in degrees
    two and three is the restriction map of inhomogeneous cochains along $f$,
    $x\mapsto x\circ(f\times\cdots\times f)$, whose coefficient map is the
    identity of $\C^\times$.
\end{lemma}

\begin{proof}\leanok
    \uses{def:mpug_h2_pullback, def:mpug_cochain_restriction,
        def:mpug_h3_group_cohomology}
    Both sides send a cochain $x$ to $(g_1,\dots,g_n)\mapsto x(f(g_1),\dots,f(g_n))$.
\end{proof}

% Partial support for #7361 (child #7754), not the full prop:BI_psi criterion.
% No U(1) restriction naturality or tensor-level block-independence claim.
\begin{lemma}[Exact extension of a subgroup cocycle]
    \label{thm:mpug_h2_exact_extension}
    \lean{TNLean.Algebra.ScalarCocycle.H2π_mem_range_map_subtype_iff,
        TNLean.Algebra.ScalarCocycle.isMulCocycle₂_coboundary_mul}
    \leanok
    \uses{def:mpug_h2_pullback, def:is_cocycle}
    Let $H\le G$ be any subgroup, let $\iota:H\hookrightarrow G$ be its
    inclusion, and let $\psi:H\times H\to\C^\times$ be a $2$-cocycle.
    Then $[\psi]$ belongs to the image of $\iota^*$ if and only if there is
    a $2$-cocycle $\Psi:G\times G\to\C^\times$ with
    $\Psi(a,b)=\psi(a,b)$ for every $a,b\in H$.
    No normalization of either cocycle is required.
    This is the representative adjustment implicit in the extension
    condition of \cite[lines~7042--7064]{FrancoRubio2025MPUGauging},
    here stated with $\C^\times$ coefficients.
    % Source: Papers/2502.20257/main.tex:7042--7064.
\end{lemma}

\begin{proof}\leanok
    \uses{def:mpug_h2_pullback, def:cohomologous_to, def:is_cocycle}
    If $[\psi]=\iota^*[\omega]$, there is a cochain
    $\varphi:H\to\C^\times$ such that
    $\psi(a,b)=\varphi(a)\varphi(b)\varphi(ab)^{-1}\omega(a,b)$.
    Extending $\varphi$ to a function $\Phi:G\to\C^\times$ that takes value $1$
    outside $H$, define
    $\Psi(g,h)=\Phi(g)\Phi(h)\Phi(gh)^{-1}\omega(g,h)$.
    In the cocycle equation for $\Psi$, cancellation leaves the common
    factor $\Phi(g)\Phi(h)\Phi(k)\Phi(ghk)^{-1}$ multiplied by the two sides
    of the cocycle equation for $\omega$. Hence $\Psi$ is a cocycle, and
    $\Phi|_H=\varphi$ gives $\Psi|_{H\times H}=\psi$ exactly.
    Conversely, any exact extension $\Psi$ satisfies
    $\iota^*[\Psi]=[\psi]$.
\end{proof}

\begin{lemma}[Circle and nonzero-complex coefficients in degree two]
    \label{thm:mpug_circle_complex_units_h2}
    \lean{TNLean.Algebra.Complex.div_toUnits_unitsPhase,
        TNLean.Algebra.ScalarCocycle.circlePhase_isMulCocycle₂,
        TNLean.Algebra.ScalarCocycle.circlePhaseInclusion_isMulCocycle₂,
        TNLean.Algebra.ScalarCocycle.H2π_circlePhaseInclusion,
        TNLean.Algebra.ScalarCocycle.circlePhaseMathlibCocycle,
        TNLean.Algebra.ScalarCocycle.H2π_map_circlePhaseMathlibCocycle,
        TNLean.Algebra.circleToComplexUnits_map_injective,
        TNLean.Algebra.circleToComplexUnits_map_surjective,
        TNLean.Algebra.circleToComplexUnits_map_bijective}
    \leanok
    \uses{thm:mpug_h2_group_cohomology_comparison,
        thm:mpug_positive_generalized_cocycle}
    Let $G$ be finite, with $U(1)$ and $\C^\times$ carrying the trivial
    $G$-actions.  Inclusion of the unit circle induces an equivalence
    \begin{align}
        H^2(G,U(1))
         & \simeq H^2(G,\C^\times).
        \label{eq:mpug_circle_complex_units_h2}
    \end{align}
    More concretely, for a $\C^\times$-valued cocycle $\omega$, define its
    pointwise phase and radial residual by
    \begin{align}
        u(g,h)
         & =\frac{\omega(g,h)}{|\omega(g,h)|}\in U(1),
         &
        r(g,h)
         & =\frac{\omega(g,h)}{u(g,h)}=|\omega(g,h)|.
        \label{eq:mpug_circle_phase_residual}
    \end{align}
    Then $u$ is a cocycle, $r$ is a positive coboundary, and hence $\omega$
    is cohomologous to the inclusion of $u$.  The inclusion map is injective
    for every group because pointwise phase is its left inverse; finiteness is
    used only to prove surjectivity.

    As recorded in \cite{gap:fbc25_circle_complex_units_cohomology}, this
    algebraic group-cohomology statement for a finite discrete group replaces
    the comparison of sheaf cohomology with topological coefficient groups
    invoked in the footnote at
    \cite[line~1931]{FrancoRubio2025MPUGauging}.
\end{lemma}

\begin{proof}\leanok
    \uses{thm:mpug_positive_generalized_cocycle}
    The phase map $q(z)=z/|z|$ is multiplicative and satisfies
    $q(u)=u$ for $u\in U(1)$, so it preserves the cocycle equation and
    retracts inclusion.  The polar identity in
    (\ref{eq:mpug_circle_phase_residual}) identifies the residual $r$ with a
    positive cocycle.  Applying
    Lemma~\ref{thm:mpug_positive_generalized_cocycle} to the one-point block
    space writes $r=d\chi$.  Therefore $\omega\sim u$, which proves
    surjectivity of coefficient inclusion on degree-two classes.  Its left
    inverse induced by $q$ proves injectivity.
\end{proof}

\subsection{Cocycle-twisted regular matrices}

\begin{definition}[Cocycle-twisted regular matrices]
    \label{def:mpug_twisted_regular_matrices}
    \lean{TNLean.Algebra.ScalarCocycle.phase,
        TNLean.Algebra.ScalarCocycle.twistedLeftRegular,
        TNLean.Algebra.ScalarCocycle.twistedRightRegular}
    \leanok
    \uses{def:scalar_cocycle, thm:mpug_circle_complex_units_h2}
    Let $G$ be a group and let $\omega$ be a $\C^\times$-valued scalar
    cochain.  Write
    \begin{align}
        \phi(g,h) & =\frac{\omega(g,h)}{|\omega(g,h)|}\in U(1).
    \end{align}
    On the standard basis of $\C^G$, define
    \begin{align}
        L_g^\phi
         & =\sum_{h\in G}\phi(g,h)^{-1}|gh\rangle\langle h|,
         &
        R_g^\phi
         & =\sum_{h\in G}\phi(h,g^{-1})^{-1}|hg^{-1}\rangle\langle h|.
        \label{eq:mpug_twisted_regular_matrices}
    \end{align}
    Whereas \cite[lines~412--418]{FrancoRubio2025MPUGauging} starts from a
    $U(1)$-valued cocycle $\Omega$, taking $\Omega=\phi$ reproduces its leg and
    inverse conventions.
\end{definition}

\begin{lemma}[Cocycle equation for the circle phase]
    \label{lem:mpug_circle_phase_cocycle}
    \lean{TNLean.Algebra.ScalarCocycle.phase_isCocycle}
    \leanok
    \uses{def:is_cocycle, thm:mpug_circle_complex_units_h2}
    If $\omega$ is a cocycle, then its circle phase satisfies
    \begin{align}
        \phi(g,h)\phi(gh,k)
         & =\phi(g,hk)\phi(h,k)
    \end{align}
    for all $g,h,k\in G$.
\end{lemma}

\begin{proof}\leanok
    \uses{thm:mpug_circle_complex_units_h2}
    The pointwise phase map $z\mapsto z/|z|$ is multiplicative, so it preserves
    the cocycle equation.
\end{proof}

\begin{lemma}[Action of the twisted regular matrices]
    \label{thm:mpug_twisted_regular_action}
    \lean{TNLean.Algebra.ScalarCocycle.twistedLeftRegular_mulVec_apply,
        TNLean.Algebra.ScalarCocycle.twistedRightRegular_mulVec_apply}
    \leanok
    \uses{def:mpug_twisted_regular_matrices}
    If $G$ is finite, then for $v\in\C^G$ and $g,h\in G$,
    \begin{align}
        (L_g^\phi v)(gh)
         & =\phi(g,h)^{-1}v(h),
         &
        (R_g^\phi v)(hg^{-1})
         & =\phi(h,g^{-1})^{-1}v(h).
    \end{align}
    % Source anchor: arXiv:2502.20257, lines 412--418.
\end{lemma}

\begin{proof}\leanok
    \uses{def:mpug_twisted_regular_matrices}
    In each matrix-vector product, exactly one summand is nonzero: the column
    indexed by $h$.
\end{proof}

\begin{lemma}[Multiplication laws for twisted regular matrices]
    \label{thm:mpug_twisted_regular_multiplication}
    \lean{TNLean.Algebra.ScalarCocycle.twistedLeftRegular_mul,
        TNLean.Algebra.ScalarCocycle.twistedRightRegular_mul}
    \leanok
    \uses{def:mpug_twisted_regular_matrices,
        lem:mpug_circle_phase_cocycle}
    If $G$ is finite and $\omega$ is a cocycle, then
    \begin{align}
        L_g^\phi L_h^\phi
         & =\phi(g,h)^{-1}L_{gh}^\phi,           \\
        R_g^\phi R_h^\phi
         & =\phi(h^{-1},g^{-1})^{-1}R_{gh}^\phi.
        \label{eq:mpug_twisted_regular_multiplication}
    \end{align}
    % Source anchor: arXiv:2502.20257, lines 406--418; the laws follow from
    % the displayed twisted regular matrices and the cocycle equation.
\end{lemma}

\begin{proof}\leanok
    \uses{lem:mpug_circle_phase_cocycle}
    Expanding the matrix products entrywise reduces the two scalar identities
    to the cocycle equation at $(g,h,k)$ and $(k,h^{-1},g^{-1})$, respectively.
\end{proof}

\begin{lemma}[Commutation of twisted left and right actions]
    \label{thm:mpug_twisted_regular_commutation}
    \lean{TNLean.Algebra.ScalarCocycle.twistedLeftRegular_commute}
    \leanok
    \uses{def:mpug_twisted_regular_matrices}
    If $G$ is finite and $\omega$ is a cocycle, then
    \begin{align}
        L_g^\phi R_h^\phi
         & =R_h^\phi L_g^\phi.
        \label{eq:mpug_twisted_regular_commutation}
    \end{align}
    % Source anchor: arXiv:2502.20257, lines 406--418 and 457; the latter
    % explicitly attributes commuting local symmetries to commutativity of the
    % left and right group actions.
\end{lemma}

\begin{proof}\leanok
    \uses{lem:mpug_circle_phase_cocycle}
    Expanding both matrix products entrywise and applying the cocycle equation
    at $(g,k,h^{-1})$ yields the equality.
\end{proof}

\begin{lemma}[Unitarity of twisted regular matrices]
    \label{thm:mpug_twisted_regular_unitary}
    \lean{TNLean.Algebra.ScalarCocycle.twistedLeftRegular_mem_unitaryGroup,
        TNLean.Algebra.ScalarCocycle.twistedRightRegular_mem_unitaryGroup}
    \leanok
    \uses{def:mpug_twisted_regular_matrices}
    If $G$ is finite, the matrices $L_g^\phi$ and $R_g^\phi$ are unitary for
    every $g\in G$.
    % Source anchor: arXiv:2502.20257, lines 406--418, where these matrices
    % form the two gauge legs of the local unitary representation.
\end{lemma}

\begin{proof}\leanok
    \uses{def:mpug_twisted_regular_matrices}
    Each matrix is a permutation matrix whose nonzero entries have modulus
    one.  Distinct columns therefore have disjoint support, while every column
    has squared norm one.
\end{proof}

\begin{lemma}[Onsite Gauss coefficient identity]
    \label{thm:mpug_onsite_gauss_phase}
    \lean{TNLean.Algebra.ScalarCocycle.onsiteGauss_phase_div}
    \leanok
    \uses{def:mpug_twisted_regular_matrices,
        lem:mpug_circle_phase_cocycle}
    Suppose that $\omega$ is a cocycle and impose the additional normalization
    \begin{align}
        \phi(t^{-1},t) & =1
    \end{align}
    for every $t\in G$.  Then
    \begin{align}
        \frac{\phi(ag^{-1},gb)}{\phi(a,b)}
         & =\frac{1}{\phi(a,g^{-1})\phi(g,b)}.
        \label{eq:mpug_onsite_gauss_phase}
    \end{align}
    The scalar convention $S_{t^{-1}}=S_t^{-1}$ appears in
    \cite[line~238]{FrancoRubio2025MPUGauging}, and the coefficient rewrite is
    at lines 3763--3766. The displayed normalization remains an explicit
    assumption here; it is not a consequence of cocyclicity.
\end{lemma}

\begin{proof}\leanok
    \uses{lem:mpug_circle_phase_cocycle}
    Applying the cocycle equation first at $(a,g^{-1},gb)$ and then at
    $(g^{-1},g,b)$, together with $\phi(e,b)=1$ from the normalization at
    $t=e$, yields the identity.
\end{proof}

\subsection{Block independence}

\begin{definition}[Block scalar and block independence]
    \label{def:mpug_block_independence_factor}
    \lean{TNLean.Algebra.LSymbol.ell,
        TNLean.Algebra.LSymbol.IsBlockIndependentEll}
    \leanok
    \uses{def:mpug_l_symbols}
    For an $L$-symbol, define
    \begin{align}
        \ell^x_{a,b;r}
         & =\frac{L^x_{ar,r^{-1}b}}{L^x_{a,b}}.
        \label{eq:mpug_block_scalar}
    \end{align}
    The scalar $\ell$ is block-independent when
    \begin{align}
        \ell^x_{a,b;r}
         & =\ell^y_{a,b;r}
        \label{eq:mpug_block_independence}
    \end{align}
    for all $a,b,r\in G$ and $x,y\in\mathsf X$.
    % Source anchor: arXiv:2502.20257, app:block_indep, lines 6965--6974.
\end{definition}

\begin{definition}[Four scalar gauge formulations]
    \label{def:mpug_block_independence_gauge_forms}
    \lean{TNLean.Algebra.LSymbol.IsBlockConstant,
        TNLean.Algebra.LSymbol.IsTrivialEll,
        TNLean.Algebra.LSymbol.IsTrivial,
        TNLean.Algebra.LSymbol.HasTrivialGauge,
        TNLean.Algebra.LSymbol.HasBlockConstantGauge,
        TNLean.Algebra.LSymbol.HasTrivialEllGauge,
        TNLean.Algebra.LSymbol.IsBlockIndependent}
    \leanok
    \uses{def:mpug_l_gauge, def:mpug_block_independence_factor}
    Writing $L'=(\beta,\gamma)\mathbin{\triangleright}L$ for a joint scalar
    gauge of an $L$-symbol, the four scalar $L$-symbol properties used in the
    proof of the source proposition are, respectively, the existence of $\beta$
    and $\gamma$ such that
    \begin{align}
        (L')^x_{g,h}       & =1
                           &                     & \text{for all $x,g,h$},
        \label{eq:mpug_block_form_trivial}                                     \\
        (L')^x_{a,b}       & =(L')^y_{a,b}
                           &                     & \text{for all $x,y,a,b$},
        \label{eq:mpug_block_form_constant}                                    \\
        \ell[L']^x_{a,b;r} & =1
                           &                     & \text{for all $x,a,b,r$},
        \label{eq:mpug_block_form_trivial_ell}                                 \\
        \ell[L']^x_{a,b;r} & =\ell[L']^y_{a,b;r}
                           &                     & \text{for all $x,y,a,b,r$}.
        \label{eq:mpug_block_form_independent_ell}
    \end{align}
    % Source anchor: arXiv:2502.20257, prop:technical01, lines 6975--6989.
\end{definition}

\begin{lemma}[Corrected block factorization]
    \label{thm:mpug_block_factorization}
    \lean{TNLean.Algebra.LSymbol.IsCompatible.apply_right_one,
        TNLean.Algebra.LSymbol.blockFactor,
        TNLean.Algebra.LSymbol.factorization_of_isBlockIndependentEll}
    \leanok
    \uses{def:mpug_l_symbols, def:mpug_block_independence_factor}
    Suppose that $L$ is compatible with $\omega$ and that $\ell$ is
    block-independent. Fix $x_0\in\mathsf X$ and set
    \begin{align}
        A(g,h)
         & =\bigl(\omega(gh,e,e)\ell^{x_0}_{g,h;h}\bigr)^{-1}.
        \label{eq:mpug_block_factor}
    \end{align}
    Then, with the reciprocal factor recorded in
    \cite{gap:fbc25_block_independence_reciprocal_errors}, for every
    $x\in\mathsf X$ and $g,h\in G$,
    \begin{align}
        L^x_{g,h}
         & =A(g,h)L^x_{e,e}.
        \label{eq:mpug_corrected_block_factorization}
    \end{align}
    % Local correction:
    % docs/paper-gaps/fbc25_block_independence_reciprocal_errors.tex.
    % Source anchor: arXiv:2502.20257, eq:Lfact, lines 6994--7002.
\end{lemma}

\begin{proof}\leanok
    \uses{def:mpug_block_independence_factor}
    Definition~\ref{def:mpug_block_independence_factor} gives
    \begin{align}
        L^x_{g,h}
         & =\frac{L^x_{gh,e}}{\ell^x_{g,h;h}}.
        \label{eq:mpug_block_factorization_first}
    \end{align}
    Block independence replaces $x$ by $x_0$ in the denominator.
    Compatibility at $(gh,e,e)$ gives
    $L^x_{gh,e}=L^x_{e,e}/\omega(gh,e,e)$, which proves
    (\ref{eq:mpug_corrected_block_factorization}).
\end{proof}

\begin{lemma}[Relative block character]
    \label{thm:mpug_relative_block_character}
    \lean{TNLean.Algebra.LSymbol.IsCompatible.apply_left_one,
        TNLean.Algebra.LSymbol.relativeCharacter,
        TNLean.Algebra.LSymbol.relativeCharacter_mul_base,
        TNLean.Algebra.LSymbol.relativeCharacter_mul}
    \leanok
    \uses{thm:mpug_block_factorization}
    Under the hypotheses of Lemma~\ref{thm:mpug_block_factorization}, define
    \begin{align}
        \rho(g)
         & =\bigl(\ell^{x_0}_{e,g;g}\omega(g,e,e)
        \omega(e,e,g)\bigr)^{-1}.
        \label{eq:mpug_relative_block_character}
    \end{align}
    As recorded in \cite{gap:fbc25_block_independence_reciprocal_errors}, the
    reciprocal factor $(\ell^{x_0}_{e,g;g})^{-1}$ corrects the corresponding
    printed factor. Then
    \begin{align}
        L^{g x}_{e,e}
         & =\rho(g)L^x_{e,e},
        \label{eq:mpug_relative_block_transport} \\
        \rho(gh)
         & =\rho(g)\rho(h).
        \label{eq:mpug_relative_block_multiplicative}
    \end{align}
    Thus $\rho\colon G\to\C^\times$ is a one-dimensional character. No
    transitivity or finiteness assumption on the action is required.
    % Local correction:
    % docs/paper-gaps/fbc25_block_independence_reciprocal_errors.tex.
    % Source anchor: arXiv:2502.20257, eq:defrho, lines 7003--7018.
\end{lemma}

\begin{proof}\leanok
    \uses{thm:mpug_block_factorization}
    Compatibility at $(e,e,g)$ gives
    $L^x_{e,g}=\omega(e,e,g)L^{g x}_{e,e}$, and comparing this with
    (\ref{eq:mpug_corrected_block_factorization}) at $(e,g)$ yields
    (\ref{eq:mpug_relative_block_transport}). Applying that identity first
    to $h$ and then to $g$, and comparing with its application to $gh$, proves
    (\ref{eq:mpug_relative_block_multiplicative}) by cancellation.
\end{proof}

\begin{lemma}[Explicit inverse gauge]
    \label{thm:mpug_block_independence_inverse_gauge}
    \lean{TNLean.Algebra.LSymbol.factorFusionGauge,
        TNLean.Algebra.LSymbol.factorActionGauge,
        TNLean.Algebra.LSymbol.gauge_one_eq_of_isBlockIndependentEll,
        TNLean.Algebra.LSymbol.gauge_inv_eq_one_of_isBlockIndependentEll,
        TNLean.Algebra.LSymbol.hasTrivialGauge_of_isBlockIndependentEll}
    \leanok
    \uses{def:mpug_l_gauge, thm:mpug_l_gauge_composition,
        thm:mpug_block_factorization, thm:mpug_relative_block_character}
    With $A$ and $\rho$ as above, define the forward gauges by
    \begin{align}
        \beta^+_{g,h}
         & =\frac{A(g,h)}{\rho(h)},
        \label{eq:mpug_forward_fusion_gauge} \\
        \gamma^{+,x}_g
         & =\bigl(L^{g x}_{e,e}\bigr)^{-1}.
        \label{eq:mpug_forward_action_gauge}
    \end{align}
    Their joint gauge multiplier equals $L$:
    \begin{align}
        (\beta^+,\gamma^+)\mathbin{\triangleright}1
         & =L.
        \label{eq:mpug_forward_gauge_multiplier}
    \end{align}
    Consequently the pointwise inverse gauges trivialize $L$:
    \begin{align}
        ((\beta^+)^{-1},(\gamma^+)^{-1})
        \mathbin{\triangleright}L
         & =1.
        \label{eq:mpug_inverse_gauge_trivialization}
    \end{align}
    As recorded in \cite{gap:fbc25_block_independence_reciprocal_errors}, the
    distinction between the forward multiplier and its pointwise inverse
    corrects the orientation of the displayed source calculation.
    If $\mathsf X$ is nonempty, a choice of $x_0$ therefore gives a
    trivializing joint scalar gauge for every compatible block-independent
    $L$-symbol.
    % Local correction:
    % docs/paper-gaps/fbc25_block_independence_reciprocal_errors.tex.
    % Source anchor: arXiv:2502.20257, eq:Lgauge and lines 7019--7027.
\end{lemma}

\begin{proof}\leanok
    \uses{thm:mpug_l_gauge_composition, thm:mpug_block_factorization,
        thm:mpug_relative_block_character}
    Substituting (\ref{eq:mpug_forward_fusion_gauge}) and
    (\ref{eq:mpug_forward_action_gauge}) into
    (\ref{eq:mpug_l_gauge}), equation
    (\ref{eq:mpug_relative_block_transport}) reduces the multiplier to
    $A(g,h)L^x_{e,e}$, which equals $L^x_{g,h}$ by
    (\ref{eq:mpug_corrected_block_factorization}). Joint gauges compose by
    pointwise multiplication, so the inverse gauges give
    (\ref{eq:mpug_inverse_gauge_trivialization}).
\end{proof}

\begin{lemma}[Trivial and block-constant gauges]
    \label{thm:mpug_trivial_iff_block_constant}
    \lean{TNLean.Algebra.LSymbol.IsTrivial.isBlockConstant,
        TNLean.Algebra.LSymbol.IsBlockConstant.isBlockIndependentEll,
        TNLean.Algebra.LSymbol.HasTrivialGauge.hasBlockConstantGauge,
        TNLean.Algebra.LSymbol.HasBlockConstantGauge.isBlockIndependent,
        TNLean.Algebra.LSymbol.hasTrivialGauge_iff_hasBlockConstantGauge}
    \leanok
    \uses{def:mpug_block_independence_gauge_forms,
        thm:mpug_l_gauge_composition, thm:mpug_l_gauge_preservation,
        thm:mpug_block_independence_inverse_gauge}
    Let $L$ be an $L$-symbol compatible with a scalar $3$-cochain,
    and suppose that $\mathsf X$ is nonempty. A joint scalar gauge trivializes
    $L$ if and only if a joint scalar gauge makes $L^x_{g,h}$ independent of
    $x$.
    % Source anchor: arXiv:2502.20257, prop:technical01, items 1 and 2.
\end{lemma}

\begin{proof}\leanok
    \uses{thm:mpug_l_gauge_composition, thm:mpug_l_gauge_preservation,
        thm:mpug_block_independence_inverse_gauge}
    A trivial $L$-symbol is block-constant. Conversely, a block-constant
    representative has block-independent $\ell$, so applying
    Lemma~\ref{thm:mpug_block_independence_inverse_gauge} to that
    representative and composing the two joint gauges trivializes $L$.
\end{proof}

\begin{lemma}[Trivial gauges and trivial block scalars]
    \label{thm:mpug_trivial_iff_trivial_ell}
    \lean{TNLean.Algebra.LSymbol.IsTrivial.isTrivialEll,
        TNLean.Algebra.LSymbol.IsTrivialEll.isBlockIndependentEll,
        TNLean.Algebra.LSymbol.HasTrivialGauge.hasTrivialEllGauge,
        TNLean.Algebra.LSymbol.HasTrivialEllGauge.isBlockIndependent,
        TNLean.Algebra.LSymbol.hasTrivialGauge_iff_hasTrivialEllGauge}
    \leanok
    \uses{def:mpug_block_independence_gauge_forms,
        thm:mpug_l_gauge_composition, thm:mpug_l_gauge_preservation,
        thm:mpug_block_independence_inverse_gauge}
    Let $L$ be an $L$-symbol compatible with a scalar $3$-cochain,
    and suppose that $\mathsf X$ is nonempty. A joint scalar gauge trivializes
    $L$ if and only if a joint scalar gauge makes every
    $\ell^x_{a,b;r}$ equal to $1$.
    % Source anchor: arXiv:2502.20257, prop:technical01, items 1 and 3.
\end{lemma}

\begin{proof}\leanok
    \uses{thm:mpug_l_gauge_composition, thm:mpug_l_gauge_preservation,
        thm:mpug_block_independence_inverse_gauge}
    A trivial $L$-symbol has $\ell=1$. Conversely, the identity $\ell=1$
    implies block independence, so applying
    Lemma~\ref{thm:mpug_block_independence_inverse_gauge} and composing the
    gauges trivializes $L$.
\end{proof}

\begin{lemma}[Scalar gauge triviality and block independence]
    \label{thm:mpug_trivial_iff_block_independent}
    \lean{TNLean.Algebra.LSymbol.IsBlockIndependent.hasTrivialGauge,
        TNLean.Algebra.LSymbol.hasTrivialGauge_iff_isBlockIndependent}
    \leanok
    \uses{def:mpug_block_independence_gauge_forms,
        thm:mpug_l_gauge_composition, thm:mpug_l_gauge_preservation,
        thm:mpug_block_independence_inverse_gauge}
    Let $L$ be an $L$-symbol compatible with a scalar $3$-cochain,
    and suppose that $\mathsf X$ is nonempty. A joint scalar gauge trivializes
    $L$ if and only if a joint scalar gauge makes $\ell^x_{a,b;r}$ independent
    of $x$.
    % Source anchor: arXiv:2502.20257, prop:technical01, items 1 and 4.
\end{lemma}

\begin{proof}\leanok
    \uses{thm:mpug_l_gauge_composition, thm:mpug_l_gauge_preservation,
        thm:mpug_block_independence_inverse_gauge}
    Gauge triviality implies block constancy and hence block independence.
    For the converse, apply
    Lemma~\ref{thm:mpug_block_independence_inverse_gauge} to the
    block-independent representative and compose its trivializing gauge with
    the original gauge.
\end{proof}

\begin{lemma}[Four scalar gauge formulations of block independence]
    \label{thm:mpug_block_independence_four_forms}
    \lean{TNLean.Algebra.LSymbol.hasTrivialGauge_tfae}
    \leanok
    \uses{def:mpug_block_independence_gauge_forms,
        thm:mpug_trivial_iff_block_constant,
        thm:mpug_trivial_iff_trivial_ell,
        thm:mpug_trivial_iff_block_independent}
    Let $L$ be an $L$-symbol compatible with a scalar $3$-cochain,
    and suppose that $\mathsf X$ is nonempty. The four scalar gauge
    properties (\ref{eq:mpug_block_form_trivial}),
    (\ref{eq:mpug_block_form_constant}),
    (\ref{eq:mpug_block_form_trivial_ell}), and
    (\ref{eq:mpug_block_form_independent_ell}) of $L$ are equivalent. This is
    the scalar core of the block-independence proposition of
    \cite{FrancoRubio2025MPUGauging}, whose fourth condition is stated for an
    MPS--MPU pair.
    % Source anchor: arXiv:2502.20257, prop:technical01, lines 6975--6989.
\end{lemma}

\begin{proof}\leanok
    \uses{thm:mpug_trivial_iff_block_constant,
        thm:mpug_trivial_iff_trivial_ell,
        thm:mpug_trivial_iff_block_independent}
    Lemmas~\ref{thm:mpug_trivial_iff_block_constant},
    \ref{thm:mpug_trivial_iff_trivial_ell}, and
    \ref{thm:mpug_trivial_iff_block_independent} identify each of the last
    three properties with the first.
\end{proof}

For a compatible scalar $L$-symbol,
Lemma~\ref{thm:mpug_block_independence_four_forms} establishes the
algebraic equivalence among its four scalar gauge properties; the corresponding
tensor-level bridge to an underlying MPS--MPU pair is not asserted here.
% Scope boundary: the tensor-facing bridge for arXiv:2502.20257,
% prop:technical01 is not contained in the scalar L-symbol declarations.

\subsection{\texorpdfstring{$L$}{L}-symbols from stabilizer cocycles}

\begin{lemma}[Stabilizer transition elements]
    \label{thm:mpug_stabilizer_transition}
    \lean{TNLean.Algebra.StabilizerRepresentatives.transitionElement_mul,
        TNLean.Algebra.StabilizerRepresentatives,
        TNLean.Algebra.StabilizerRepresentatives.ofPretransitive,
        TNLean.Algebra.StabilizerRepresentatives.transitionElement,
        TNLean.Algebra.StabilizerRepresentatives.coe_transitionElement,
        TNLean.Algebra.StabilizerRepresentatives.transitionElement_stabilizer}
    \leanok
    Suppose that $G$ acts transitively on $\mathsf X$, and fix
    $x_0\in\mathsf X$. There are representatives $k_x\in G$ satisfying
    $k_xx_0=x$ and $k_{x_0}=e$. With
    $H=\operatorname{Stab}_G(x_0)$, define
    \begin{align}
        t_k(g,x)
         & =k_{g x}^{-1}gk_x.
        \label{eq:mpug_stabilizer_transition}
    \end{align}
    Then $t_k(g,x)\in H$ and
    \begin{align}
        t_k(g_1g_2,x)
         & =t_k(g_1,g_2x)t_k(g_2,x).
        \label{eq:mpug_stabilizer_transition_product}
    \end{align}
    Consequently, the source factors are
    \begin{align}
        h_2
         & =t_k(g_2,x)=k_{g_2x}^{-1}g_2k_x,
        \label{eq:mpug_stabilizer_h2}                  \\
        h_1
         & =t_k(g_1,g_2x)=k_{g_1g_2x}^{-1}g_1k_{g_2x}.
        \label{eq:mpug_stabilizer_h1}
    \end{align}
    As recorded in \cite{gap:fbc25_stabilizer_transition_index_typo}, the last
    representative in $h_1$ corrects the printed $k_{g_1x}$. No finiteness or
    normality assumption is required.
    % Local correction:
    % docs/paper-gaps/fbc25_stabilizer_transition_index_typo.tex.
    % Source anchor: arXiv:2502.20257, lemma:h1h2 and eq:h1h2,
    % lines 7028--7042.
\end{lemma}

\begin{proof}\leanok
    Transitivity gives one element carrying $x_0$ to each $x$, with the
    identity chosen at $x=x_0$. Acting with (\ref{eq:mpug_stabilizer_transition})
    on $x_0$ gives
    $k_{gx}^{-1}g(k_xx_0)=k_{gx}^{-1}(gx)=x_0$.
    For (\ref{eq:mpug_stabilizer_transition_product}), substituting
    (\ref{eq:mpug_stabilizer_transition}) and cancelling
    $k_{g_2x}k_{g_2x}^{-1}$ yields the product formula. The two formulas for
    $h_1$ and $h_2$ are the corresponding specializations.
\end{proof}

\begin{lemma}[L-symbol induced by a stabilizer cocycle]
    \label{thm:mpug_stabilizer_cocycle_l_symbol}
    \lean{TNLean.Algebra.StabilizerRepresentatives.inducedLSymbol_isCompatible,
        TNLean.Algebra.StabilizerRepresentatives.inducedLSymbol,
        TNLean.Algebra.StabilizerRepresentatives.inducedLSymbol_base,
        TNLean.Algebra.StabilizerRepresentatives.inducedActionGauge,
        TNLean.Algebra.StabilizerRepresentatives.inducedLSymbol_eq_gauge,
        TNLean.Algebra.StabilizerRepresentatives.isMulCoboundary₂_div_iff_exists_actionGauge_inducedLSymbol_eq}
    \leanok
    \uses{def:mpug_l_symbols, def:mpug_l_gauge, def:is_cocycle,
        def:cohomologous_to, thm:mpug_stabilizer_transition}
    Let $H=\operatorname{Stab}_G(x_0)$ and let
    $\psi\colon H\times H\to\C^\times$ be a scalar $2$-cocycle. Define the
    $\omega=1$ scalar specialization of \cite[Eq.~(20)]{GarreRubio2023MPOSym}
    by
    \begin{align}
        L^x_{g,h}
         & =\psi\bigl(t_k(g,hx),t_k(h,x)\bigr).
        \label{eq:mpug_induced_stabilizer_l_symbol}
    \end{align}
    The $L$-symbol (\ref{eq:mpug_induced_stabilizer_l_symbol}) satisfies
    (\ref{eq:mpug_l_compatibility}) with $\omega=1$, and its base-point
    restriction recovers $\psi$:
    \begin{align}
        L^{x_0}_{a,b} & =\psi(a,b)
        \label{eq:mpug_induced_stabilizer_base}
    \end{align}
    for all $a,b\in H$.

    More precisely, two stabilizer cocycles are cohomologous exactly when
    their induced $L$-symbols differ by an action-tensor gauge and the trivial
    fusion-tensor gauge:
    \begin{align}
        \psi_1\sim\psi_2
         & \quad\Longleftrightarrow\quad
        \exists\gamma,\quad
        L[\psi_1]=(1,\gamma)\mathbin{\triangleright}L[\psi_2].
        \label{eq:mpug_induced_cohomology_iff}
    \end{align}
    If the cohomology relation is given by
    $\varphi\colon H\to\C^\times$, with
    \begin{align}
        \psi_1(a,b)
         & =\varphi(a)\varphi(b)\varphi(ab)^{-1}\psi_2(a,b),
        \label{eq:mpug_stabilizer_cocycle_gauge}
    \end{align}
    then the forward action-tensor gauge is
    \begin{align}
        \gamma_{g,x} & =\varphi\bigl(t_k(g,x)\bigr)^{-1}.
        \label{eq:mpug_induced_action_gauge}
    \end{align}
    Thus the induced action-gauge class depends only on the cohomology class
    of $\psi$.

    This establishes the scalar part of the $L$-symbol reduction in
    \cite[lines~7028--7042]{FrancoRubio2025MPUGauging}; the
    representation-theoretic claim for a tensor-derived $L$-symbol is not
    asserted here.
    % Source anchor: arXiv:2203.12563, Eq. (20).
\end{lemma}

\begin{proof}\leanok
    \uses{def:mpug_l_gauge, thm:mpug_stabilizer_transition}
    With
    $a=t_k(g,hkx)$, $b=t_k(h,kx)$, and $c=t_k(k,x)$,
    the product law (\ref{eq:mpug_stabilizer_transition_product}) gives
    $t_k(gh,kx)=ab$ and $t_k(hk,x)=bc$. Compatibility is therefore
    \begin{align}
        \psi(a,b)\psi(ab,c)
         & =\psi(a,bc)\psi(b,c),
    \end{align}
    which is the cocycle equation. If $p,q\in H$, then
    $t_k(p,x_0)=p$ and $t_k(q,x_0)=q$ because $k_{x_0}=e$, proving
    (\ref{eq:mpug_induced_stabilizer_base}). For the forward gauge
    calculation, setting $a=t_k(g,hx)$ and $b=t_k(h,x)$ gives
    $t_k(gh,x)=ab$. Substituting (\ref{eq:mpug_induced_action_gauge}) into
    (\ref{eq:mpug_l_gauge}) gives the multiplier
    \begin{align}
        \frac{\gamma_{gh,x}}{\gamma_{g,hx}\gamma_{h,x}}
         & =\varphi(ab)^{-1}
        \bigl(\varphi(a)^{-1}\varphi(b)^{-1}\bigr)^{-1}
        =\varphi(a)\varphi(b)\varphi(ab)^{-1}.
    \end{align}
    Formula (\ref{eq:mpug_stabilizer_cocycle_gauge}) now identifies the two
    induced $L$-symbols. Conversely, if
    $L[\psi_1]=(1,\gamma)\mathbin{\triangleright}L[\psi_2]$ and
    $\varphi(a)=\gamma_{a,x_0}^{-1}$ for $a\in H$, evaluating the gauge
    equality at $x_0$ and $a,b\in H$ gives
    \begin{align}
        \psi_1(a,b)
         & =\frac{\gamma_{ab,x_0}}
        {\gamma_{a,x_0}\gamma_{b,x_0}}\psi_2(a,b)
        =\varphi(a)\varphi(b)\varphi(ab)^{-1}\psi_2(a,b).
    \end{align}
    Hence $\psi_1\sim\psi_2$, which proves the reverse implication in
    (\ref{eq:mpug_induced_cohomology_iff}).
\end{proof}

\begin{lemma}[Explicit gauge from a stabilizer cocycle extension]
    \label{thm:mpug_stabilizer_extension_gauge}
    \lean{TNLean.Algebra.StabilizerRepresentatives.extensionActionGauge,
        TNLean.Algebra.StabilizerRepresentatives.inducedLSymbol_eq_extensionActionGauge,
        TNLean.Algebra.StabilizerRepresentatives.inducedLSymbol_hasTrivialGauge_of_extension}
    \leanok
    \uses{thm:mpug_stabilizer_cocycle_l_symbol, def:mpug_l_gauge,
        def:mpug_block_independence_gauge_forms}
    Let $G$ act on $\mathsf X$, with basepoint $x_0$, stabilizer
    $H=\operatorname{Stab}_G(x_0)$, and representatives $k_x$ satisfying
    $k_xx_0=x$ and $k_{x_0}=e$.
    Let $\psi\colon H\times H\to\C^\times$ and let
    $\Psi\colon G\times G\to\C^\times$ be a two-cocycle whose restriction
    to $H\times H$ is $\psi$. Define
    \begin{align}
        \gamma_g^x
         & =\frac{\Psi(k_{gx}^{-1},gk_x)\Psi(g,k_x)}
        {\Psi(e,e)\Psi(k_x,k_x^{-1})}.
        \label{eq:mpug_extension_gamma}
    \end{align}
    With $h_1=t_k(g_1,g_2x)$ and $h_2=t_k(g_2,x)$, the induced symbols satisfy
    \begin{align}
        L^x_{g_1,g_2}=\psi(h_1,h_2)
         & =\Psi(g_1,g_2)\frac{\gamma_{g_1g_2}^x}
        {\gamma_{g_1}^{g_2x}\gamma_{g_2}^x}.
        \label{eq:mpug_extension_identity}
    \end{align}
    Consequently the inverse fusion and action gauges
    $(\Psi^{-1},\gamma^{-1})$ make $L$ identically one.
    This is the scalar representative calculation in the forward proof of
    the block-independence extension proposition of
    \cite{FrancoRubio2025MPUGauging};
    here extension means the equality $\Psi|_{H\times H}=\psi$ of chosen
    cocycle representatives, rather than only equality of their classes.
    See \cite{gap:fbc25_mpu_action_stabilizer_source_audit} for the
    relationship between cocycle representatives, cohomology classes, and
    tensor data.
\end{lemma}

\begin{proof}\leanok
    With $a=k_{g_1g_2x}$, $b=k_{g_2x}$, $c=k_x$, $u=g_1$, and $v=g_2$,
    the cocycle equation gives $\Psi(z,e)=\Psi(e,z)=\Psi(e,e)$ and
    $\Psi(b^{-1},b)=\Psi(b,b^{-1})$. Applying it successively to the triples
    \begin{align}
         & (a^{-1}ub,b^{-1},vc),\quad (a^{-1}ub,b^{-1},b),\quad
        (a^{-1},u,vc),\quad (a^{-1},u,b),\quad (u,v,c)
    \end{align}
    and solving these equalities for the corresponding factors of
    $\Psi(a^{-1}ub,b^{-1}vc)$ yields
    (\ref{eq:mpug_extension_identity}) with
    (\ref{eq:mpug_extension_gamma}) after cancellation. The restriction
    hypothesis identifies its left side with $\psi(h_1,h_2)$. Substitution into
    (\ref{eq:mpug_l_gauge}) cancels the multiplier under the inverse gauges.
\end{proof}

\begin{lemma}[Exact extension from trivializing scalar gauges]
    \label{lem:mpug_trivializing_extension}
    \lean{TNLean.Algebra.LSymbol.isMulCocycle₂_of_isTrivial_gauge,
        TNLean.Algebra.LSymbol.basepointExtension,
        TNLean.Algebra.LSymbol.basepointExtension_isMulCocycle₂,
        TNLean.Algebra.StabilizerRepresentatives.basepointExtension_restrict,
        TNLean.Algebra.StabilizerRepresentatives.exists_extension_of_inducedLSymbol_hasTrivialGauge}
    \leanok
    \uses{thm:mpug_stabilizer_cocycle_l_symbol, def:mpug_l_gauge}
    Let $G$ act on $\mathsf X$ with basepoint $x_0$, stabilizer
    $H=\operatorname{Stab}_G(x_0)$, and representatives $k_x$ satisfying
    $k_xx_0=x$ and $k_{x_0}=e$. Let $\psi:H\times H\to\C^\times$ be a
    two-cocycle and $L[\psi]$ its induced scalar $L$-symbol.
    Suppose scalar cochains $\beta:G\times G\to\C^\times$ and
    $\gamma:G\times\mathsf X\to\C^\times$ satisfy
    $(\beta,\gamma)\mathbin{\triangleright}L[\psi]=1$.
    Then $\beta$ is a two-cocycle and
    \begin{align}
        \Psi(g,h)
         & =\beta(g,h)^{-1}
        \frac{\gamma_{g,x_0}\gamma_{h,x_0}}{\gamma_{gh,x_0}}
        \label{eq:mpug_trivializing_extension}
    \end{align}
    is a two-cocycle with $\Psi|_{H\times H}=\psi$ exactly.
    Neither cocycles nor gauges need be normalized.
    This is the converse scalar calculation in
    \cite[lines~7057--7064]{FrancoRubio2025MPUGauging}, with gauges that
    trivialize $L[\psi]$, the inverses of the source's forward gauges.
\end{lemma}

\begin{proof}\leanok
    \uses{thm:mpug_stabilizer_cocycle_l_symbol, def:mpug_l_gauge,
        thm:mpug_h2_exact_extension}
    Compatibility of $L[\psi]$ with $\omega=1$ transports under
    $(\beta,\gamma)$ to compatibility of the constant symbol $1$ with
    $d\beta$. Thus $d\beta=1$, which is the two-cocycle equation for
    $\beta$. Setting $\varphi(g)=\gamma_{g,x_0}$, formula
    (\ref{eq:mpug_trivializing_extension}) is
    $\Psi=d\varphi\,\beta^{-1}$, hence is a two-cocycle.
    For $a,b\in H$, the equality $bx_0=x_0$ and the chosen
    $k_{x_0}=e$ identify $L[\psi]^{x_0}_{a,b}$ with $\psi(a,b)$.
    Evaluating the trivializing-gauge identity there gives
    $\Psi(a,b)=\psi(a,b)$.
\end{proof}

\begin{lemma}[Block independence extends the stabilizer cocycle]
    \label{lem:mpug_bi_basepoint_extension}
    \lean{TNLean.Algebra.LSymbol.basepointExtension_apply_of_isTrivial_gauge,
        TNLean.Algebra.LSymbol.exists_extension_of_isBlockIndependent}
    \leanok
    \uses{lem:mpug_trivializing_extension, thm:mpug_block_independence_four_forms,
        def:mpug_block_independence_gauge_forms}
    Let $G$ act on $\mathsf X$ with basepoint $x_0$ and stabilizer
    $H=\operatorname{Stab}_G(x_0)$, and let $L$ be any $L$-symbol compatible
    with $\omega=1$. If $L$ is block independent after a joint scalar gauge,
    then there is a two-cocycle $\Psi:G\times G\to\C^\times$ with
    $\Psi(a,b)=L^{x_0}_{a,b}$ for all $a,b\in H$. This is the converse
    direction of \cite[lines~7057--7064]{FrancoRubio2025MPUGauging} for an
    arbitrary compatible $L$-symbol, not only for one induced by
    representatives.
\end{lemma}

\begin{proof}\leanok
    \uses{lem:mpug_trivializing_extension, thm:mpug_block_independence_four_forms}
    By Lemma~\ref{thm:mpug_block_independence_four_forms}, some joint scalar
    gauge $(\beta,\gamma)$ trivializes $L$. The cochain
    (\ref{eq:mpug_trivializing_extension}) is a two-cocycle by
    Lemma~\ref{lem:mpug_trivializing_extension}. For $b\in H$ the trivializing
    identity at $(x_0,a,b)$ uses only $bx_0=x_0$, and it gives
    $\Psi(a,b)=L^{x_0}_{a,b}$.
\end{proof}

% Scalar criterion; the circle-coefficient and tensor forms follow below.
\begin{lemma}[Scalar block independence and restriction of the stabilizer class]
    \label{lem:mpug_scalar_bi_restriction}
    \lean{TNLean.Algebra.StabilizerRepresentatives.inducedLSymbol_isBlockIndependent_iff_exists_extension,
        TNLean.Algebra.StabilizerRepresentatives.inducedLSymbol_isBlockIndependent_iff_H2π_mem_range_map_subtype}
    \leanok
    \uses{lem:mpug_trivializing_extension, thm:mpug_stabilizer_extension_gauge,
        thm:mpug_h2_exact_extension, thm:mpug_block_independence_four_forms}
    Let $G$ act on $\mathsf X$ with basepoint $x_0$, stabilizer
    $H=\operatorname{Stab}_G(x_0)$, and representatives $k_x$ satisfying
    $k_xx_0=x$ and $k_{x_0}=e$. For a two-cocycle
    $\psi:H\times H\to\C^\times$, the following are equivalent:
    \begin{enumerate}
        \item The induced scalar $L$-symbol $L[\psi]$ is block independent
              after a joint scalar gauge.
        \item There is a two-cocycle $\Psi:G\times G\to\C^\times$ with
              $\Psi|_{H\times H}=\psi$ exactly.
        \item The class $[\psi]$ belongs to the image of the restriction map
              $H^2(G,\C^\times)\to H^2(H,\C^\times)$ for trivial coefficients.
    \end{enumerate}
    No normality or finiteness assumption is imposed on $H$ or $G$.
    This states the scalar consequence of
    \cite[lines~7042--7064]{FrancoRubio2025MPUGauging} for the induced symbols.
    The distinction between classes, representatives, and tensor data is
    explained in \cite{gap:fbc25_mpu_action_stabilizer_source_audit}.
\end{lemma}

\begin{proof}\leanok
    \uses{thm:mpug_block_independence_four_forms,
        lem:mpug_trivializing_extension, thm:mpug_stabilizer_extension_gauge,
        thm:mpug_h2_exact_extension}
    By Lemma~\ref{thm:mpug_block_independence_four_forms}, scalar block
    independence is equivalent to the existence of a trivializing scalar
    gauge. Lemma~\ref{lem:mpug_trivializing_extension} constructs an exact
    extension from such a gauge, and
    Lemma~\ref{thm:mpug_stabilizer_extension_gauge} constructs a trivializing
    gauge from an exact extension. Finally,
    Lemma~\ref{thm:mpug_h2_exact_extension} identifies exact extension of
    $\psi$ with membership of $[\psi]$ in the image of restriction.
\end{proof}

\begin{lemma}[Reconstruction of stabilizer L-symbols]
    \label{thm:mpug_stabilizer_reconstruction}
    \lean{TNLean.Algebra.StabilizerRepresentatives.reconstructedLSymbol,
        TNLean.Algebra.StabilizerRepresentatives.reconstructedLSymbol_isCompatible,
        TNLean.Algebra.StabilizerRepresentatives.reconstructedLSymbol_one_isCompatible,
        TNLean.Algebra.StabilizerRepresentatives.reconstructedLSymbol_eq_mul,
        TNLean.Algebra.StabilizerRepresentatives.reconstructedLSymbol_base,
        TNLean.Algebra.StabilizerRepresentatives.gauge_reconstructedLSymbol_isCompatible,
        TNLean.Algebra.ScalarThreeCochain.exists_fusionGauge_eq_one_of_isTrivialGaugeClass_comap,
        TNLean.Algebra.LSymbol.IsCompatible.mul}
    \leanok
    \uses{def:mpug_l_symbols, def:mpug_l_gauge, def:mpug_three_cocycle,
        def:mpug_fusion_gauge, def:mpug_trivial_gauge_class, def:is_cocycle,
        thm:mpug_stabilizer_transition}
    Let $G$ act on $\mathsf X$ with basepoint $x_0$, stabilizer
    $H=\operatorname{Stab}_G(x_0)$, and representatives $k_x$ as in
    Lemma~\ref{thm:mpug_stabilizer_transition}. Let $\omega$ be a scalar
    $3$-cocycle on $G$ with $\omega(a,b,c)=1$ for all $a,b,c\in H$, and let
    $\psi$ be a scalar $2$-cocycle on $H$. With $h_2=t_k(g_2,x)$ and
    $h_1=t_k(g_1,g_2x)$, define
    \begin{align}
        L^x_{g_1,g_2}
         & =\frac{\omega(g_1g_2k_x,h_2^{-1},h_1^{-1})\,\psi(h_1,h_2)}
        {\omega(g_1,g_2k_x,h_2^{-1})\,\omega(g_1,g_2,k_x)}.
        \label{eq:mpug_stabilizer_reconstruction}
    \end{align}
    Then the $L$-symbol (\ref{eq:mpug_stabilizer_reconstruction}) of
    \cite[Eq.~(20)]{GarreRubio2023MPOSym} satisfies
    (\ref{eq:mpug_l_compatibility}), with
    $L^{x_0}_{a,b}=\psi(a,b)$ for $a,b\in H$.

    If only $\omega|_H$ is a coboundary, some fusion gauge $\beta$ on $G$
    makes $\beta\mathbin{\triangleright}\omega$ equal to one on
    $H\times H\times H$, and applying the inverse gauge
    $(\beta^{-1},1)$ to (\ref{eq:mpug_stabilizer_reconstruction}) built from
    $\beta\mathbin{\triangleright}\omega$ gives an $L$-symbol compatible with
    $\omega$. The hypothesis that $\omega$ is one on the stabilizer fixes the
    fusion gauge in which the source's formula holds; as recorded in
    \cite{gap:glm23_eq20_fusion_gauge}, for a general
    $\omega$ with $\omega|_H=d\beta$ the printed formula can fail.
    No normalization of $\omega$ is required.
    % Local fix: docs/paper-gaps/glm23_eq20_fusion_gauge.tex.
    % Source anchor: arXiv:2203.12563, REsubmission.tex lines 740--761.
\end{lemma}

\begin{proof}\leanok
    \uses{thm:mpug_stabilizer_transition, thm:mpug_stabilizer_cocycle_l_symbol}
    The formula factors as its value at $\psi=1$ times the induced
    $L$-symbol (\ref{eq:mpug_induced_stabilizer_l_symbol}), and pointwise
    products of compatible $L$-symbols are compatible with the product of
    the $3$-cochains. By Lemma~\ref{thm:mpug_stabilizer_cocycle_l_symbol}
    it suffices to treat $\psi=1$. For fixed $x$ and $g,h,k$, with $p=k_x$,
    $a=t_k(k,x)$, $b=t_k(h,kx)$, and $c=t_k(g,hkx)$, we have $k_{kx}=kpa^{-1}$
    and, by (\ref{eq:mpug_stabilizer_transition_product}),
    $t_k(hk,x)=ba$ and $t_k(gh,kx)=cb$. The compatibility equation then
    follows from the $3$-cocycle equation at the four quadruples
    \begin{align}
         & (g,h,k,p),\quad (g,h,kp,a^{-1}),\quad (g,hkp,a^{-1},b^{-1}),\quad
        (ghkp,a^{-1},b^{-1},c^{-1}),
    \end{align}
    together with $\omega(a^{-1},b^{-1},c^{-1})=1$. At $x_0$ and $a,b\in H$
    every argument of $\omega$ in (\ref{eq:mpug_stabilizer_reconstruction})
    lies in $H$, which gives $L^{x_0}_{a,b}=\psi(a,b)$. If
    $\omega|_H=d\beta_0$, extending $\beta_0^{-1}$ by one outside
    $H\times H$ produces a coboundary that cancels $\omega$ on
    $H\times H\times H$. Finally, $(\beta^{-1},1)$ changes
    $\beta\mathbin{\triangleright}\omega$ back to $\omega$.
\end{proof}

\begin{lemma}[Classification and torsor of stabilizer L-symbols]
    \label{thm:mpug_stabilizer_torsor}
    \lean{TNLean.Algebra.StabilizerRepresentatives.exists_actionGauge_eq_inducedLSymbol,
        TNLean.Algebra.StabilizerRepresentatives.exists_actionGauge_eq_reconstructedLSymbol,
        TNLean.Algebra.StabilizerRepresentatives.exists_actionGauge_gauge_eq_reconstructedLSymbol,
        TNLean.Algebra.StabilizerRepresentatives.isCompatible_iff_exists_actionGauge_mul_inducedLSymbol,
        TNLean.Algebra.StabilizerRepresentatives.actionGaugeEquiv_mul_inducedLSymbol_iff,
        TNLean.Algebra.StabilizerRepresentatives.solutionSetoid,
        TNLean.Algebra.StabilizerRepresentatives.h2EquivSolutionClasses,
        TNLean.Algebra.StabilizerRepresentatives.h2EquivSolutionClasses_H2π,
        TNLean.Algebra.StabilizerRepresentatives.h2EquivSolutionClasses_apply,
        TNLean.Algebra.StabilizerRepresentatives.solutionAddAction,
        TNLean.Algebra.StabilizerRepresentatives.existsUnique_vadd_eq,
        TNLean.Algebra.LSymbol.ActionGaugeEquiv,
        TNLean.Algebra.LSymbol.actionGaugeEquiv_equivalence,
        TNLean.Algebra.LSymbol.IsCompatible.inv,
        TNLean.Algebra.LSymbol.IsCompatible.mul_inv_one,
        TNLean.Algebra.LSymbol.IsCompatible.gauge_one,
        TNLean.Algebra.LSymbol.IsCompatible.isCocycle_restrict,
        TNLean.Algebra.LSymbol.gauge_one_mul}
    \leanok
    \uses{thm:mpug_stabilizer_reconstruction, def:mpug_l_gauge,
        thm:mpug_h2_group_cohomology_comparison, def:cohomologous_to}
    Keep the setting of Lemma~\ref{thm:mpug_stabilizer_reconstruction},
    and call two $L$-symbols action-gauge equivalent when they differ by
    $(1,\gamma)\mathbin{\triangleright}$ for some action-tensor gauge
    $\gamma$.
    \begin{enumerate}
        \item If $R$ is compatible with the trivial $3$-cochain, then
              $R=(1,\gamma)\mathbin{\triangleright}L[\psi]$ for the induced
              $L$-symbol (\ref{eq:mpug_induced_stabilizer_l_symbol}) of the
              $2$-cocycle $\psi(a,b)=R^{x_0}_{a,b}$, with
              $\gamma^x_g=R^{x_0}_{k_{gx},t_k(g,x)}/R^{x_0}_{g,k_x}$.
        \item If $\omega$ is a $3$-cocycle equal to one on
              $H\times H\times H$, every $L$-symbol compatible with $\omega$
              is action-gauge equivalent to
              (\ref{eq:mpug_stabilizer_reconstruction}) for some $2$-cocycle
              $\psi$ on $H$. For a general fusion gauge $\beta$ making
              $\beta\mathbin{\triangleright}\omega$ trivial on $H$, the
              same holds for $(\beta,1)\mathbin{\triangleright}L$ and the
              formula built from $\beta\mathbin{\triangleright}\omega$.
        \item Let $L_0$ be compatible with an arbitrary scalar $3$-cochain
              $\omega$. An $L$-symbol $L$ is compatible with $\omega$ if and
              only if $L$ is action-gauge equivalent to $L_0L[\psi]$ for
              some $2$-cocycle $\psi$ on $H$, and $L_0L[\psi_1]$ and
              $L_0L[\psi_2]$ are action-gauge equivalent if and only if
              $\psi_1\sim\psi_2$. Hence sending the class of $\psi$ to the
              class of $L_0L[\psi]$ is a bijection from $H^2(H,\C^\times)$
              onto the action-gauge classes of $L$-symbols compatible with
              $\omega$.
        \item Sending the class of $\psi$ and the class of $L$ to the
              class of $L\,L[\psi]$ is a well-defined action of the group
              cohomology $H^2(H,\C^\times)$, for the trivial action of $H$
              on $\C^\times$, on these classes, and for any two classes
              exactly one cohomology class carries the first to the second.
    \end{enumerate}
    The last two items are the torsor statement of \cite{GarreRubio2023MPOSym}
    for fixed $(G,\omega,H)$, with the fusion gauge held fixed as in the
    source. In all four items $G$, and hence $H$, is a small group: its
    underlying set lies in a fixed set-theoretic universe containing $\Z$,
    in which the group cohomology $H^2(H,\C^\times)$ is formed. In the last
    two items the set $\mathsf X$ lies in the same universe.
    % Source anchor: arXiv:2203.12563, REsubmission.tex lines 749--761.
\end{lemma}

\begin{proof}\leanok
    \uses{thm:mpug_stabilizer_reconstruction,
        thm:mpug_stabilizer_cocycle_l_symbol}
    For the first item, compatibility at $x_0$ with arguments
    $(g,h,k_x)$ expresses $R^x_{g,h}$ through basepoint values. Rewriting
    $hk_x=k_{hx}t_k(h,x)$ and $gk_{hx}=k_{ghx}t_k(g,hx)$ and using
    compatibility at $x_0$ with arguments $(g,k_{hx},t_k(h,x))$ and
    $(k_{ghx},t_k(g,hx),t_k(h,x))$ combines the three equations into
    $R^x_{g,h}=\gamma^x_{gh}(\gamma^{hx}_g\gamma^x_h)^{-1}
        \psi(t_k(g,hx),t_k(h,x))$. Compatibility at $x_0$ with arguments in
    $H$ is the cocycle equation for $\psi$.

    For the third item, if $L$ and $L_0$ are compatible with $\omega$,
    then $LL_0^{-1}$ is compatible with the trivial $3$-cochain, and the first
    item applies; conversely $L_0L[\psi]$ is compatible with $\omega$, and an
    action gauge with trivial fusion part preserves compatibility and
    commutes with multiplication by $L_0$. Cancelling $L_0$ reduces the
    freeness statement to (\ref{eq:mpug_induced_cohomology_iff}). The second
    item is the third with $L_0$ the formula at $\psi=1$, which is
    compatible by Lemma~\ref{thm:mpug_stabilizer_reconstruction}.
    For the last item, cocycles modulo cohomology form the second group
    cohomology of $H$ with coefficients in $\C^\times$, and
    \begin{align}
        [\psi_1]+[\psi_2] & =[\psi_1\psi_2],
        \notag                                     \\
        L[\psi_1\psi_2]   & =L[\psi_1]\,L[\psi_2].
        \notag
    \end{align}
    Cohomologous cocycles give induced $L$-symbols related by an action
    gauge, so the rule is well defined, and by the two identities acting by
    $[\psi_1]+[\psi_2]$ is acting by $[\psi_2]$ and then by $[\psi_1]$. With
    base point $L$ the action is the bijection of the third item.
\end{proof}

\begin{lemma}[Scalar block independence with circle coefficients]
    \label{lem:mpug_scalar_bi_restriction_circle}
    \lean{TNLean.Algebra.circleRepresentationRes,
        TNLean.Algebra.circleToComplexUnits_map_comp_map_res,
        TNLean.Algebra.circleToComplexUnits_map_res_apply,
        TNLean.Algebra.mem_range_map_circleRepresentationRes_iff,
        TNLean.Algebra.LSymbol.isBlockIndependent_gauge_iff,
        TNLean.Algebra.StabilizerRepresentatives.inducedLSymbol_isBlockIndependent_iff_circle}
    \leanok
    \uses{lem:mpug_scalar_bi_restriction, thm:mpug_circle_complex_units_h2}
    Restriction along a group homomorphism $f\colon S\to G$ commutes with
    the maps $H^2(-,U(1))\to H^2(-,\C^\times)$ induced by inclusion of the
    unit circle. If $G$ is finite, a class $c\in H^2(S,U(1))$ is the
    restriction of a class of $G$ exactly when its image in
    $H^2(S,\C^\times)$ is. Block independence of an $L$-symbol after
    a joint scalar gauge is unchanged by any further joint scalar gauge.

    Consequently, in the setting of Lemma~\ref{lem:mpug_scalar_bi_restriction}
    with $G$ finite, the induced $L$-symbol $L[\psi]$ is block independent
    after a joint scalar gauge if and only if the class in $H^2(H,U(1))$ of
    the pointwise phase $\psi/|\psi|$ lies in the image of the restriction map
    $H^2(G,U(1))\to H^2(H,U(1))$. For a cocycle of modulus one this phase is
    $\psi$ itself. This is the coefficient group of
    \cite[lines~7042--7064]{FrancoRubio2025MPUGauging}; the comparison of
    coefficients is the algebraic one of
    Lemma~\ref{thm:mpug_circle_complex_units_h2}.
\end{lemma}

\begin{proof}\leanok
    \uses{lem:mpug_scalar_bi_restriction, thm:mpug_circle_complex_units_h2,
        thm:mpug_l_gauge_composition}
    Both composites of restriction and inclusion are induced by the same
    homomorphism $S\to G$ and the same coefficient map
    $U(1)\to\C^\times$, so they agree. If $c=\operatorname{res}(a)$, then
    the image of $c$ is the restriction of the image of $a$. Conversely, if
    the image of $c$ is $\operatorname{res}(b)$, surjectivity of inclusion
    for the finite group $G$ writes $b$ as the image of some $a$; then $c$
    and $\operatorname{res}(a)$ have the same image, and inclusion is
    injective for every group. Gauge invariance follows from
    (\ref{eq:mpug_l_gauge_composition}) applied with the inverse gauge. By
    Lemma~\ref{thm:mpug_circle_complex_units_h2} the class of the phase of
    $\psi$ maps to $[\psi]$, and Lemma~\ref{lem:mpug_scalar_bi_restriction}
    finishes the proof.
\end{proof}

\begin{lemma}[Stabilizer form of the L-symbols of permuted blocks]
    \label{lem:mpug_permuted_block_stabilizer_form}
    \lean{MPOTensor.GroupFamily.BlockActionData.exists_rescale_lSymbol_eq_inducedLSymbol,
        MPOTensor.GroupFamily.BlockActionData.exists_rescale_lSymbol_eq_inducedLSymbol_of_finite}
    \leanok
    \uses{lem:mpug_permuted_block_l_gauge,
        thm:mpug_permuted_block_l_compatibility, cor:mpug_anomaly_class,
        thm:mpug_stabilizer_torsor, thm:mpug_stabilizer_cocycle_l_symbol,
        thm:mpug_stabilizer_transition, thm:mpug_circle_complex_units_h2}
    In the setting of Lemma~\ref{lem:mpug_permuted_block_l_gauge}, choose
    $x_0\in\mathsf X$, let $H$ be its stabilizer, and choose $k_x\in G$ with
    $k_xx_0=x$ and $k_{x_0}=e$. If the anomaly class vanishes, some rescaling
    of the fusion and action tensors gives $\omega=1$ and
    \begin{align}
        L^x_{g_1,g_2}
         & =\psi(h_1,h_2),
         &
        h_1
         & =k_{g_1g_2x}^{-1}g_1k_{g_2x},
         &
        h_2
         & =k_{g_2x}^{-1}g_2k_x,
        \notag
    \end{align}
    for a two-cocycle $\psi\colon H\times H\to\C^\times$. If $G$ is finite,
    $\psi$ can be chosen with values of modulus one.

    This is \cite[lines~7028--7042]{FrancoRubio2025MPUGauging}, with the
    corrected last representative in $h_1$ of
    \cite{gap:fbc25_stabilizer_transition_index_typo}; $\psi$ is a cocycle
    representative, not a cohomology class.
\end{lemma}

\begin{proof}\leanok
    \uses{lem:mpug_permuted_block_l_gauge,
        thm:mpug_permuted_block_l_compatibility, cor:mpug_anomaly_class,
        thm:mpug_stabilizer_torsor, thm:mpug_stabilizer_cocycle_l_symbol,
        thm:mpug_circle_complex_units_h2}
    By Corollary~\ref{cor:mpug_anomaly_class} there is $\beta$ with
    $\beta\mathbin{\triangleright}\omega=1$, and by
    Lemma~\ref{lem:mpug_permuted_block_l_gauge} rescaling the fusion tensors
    by $\beta$ makes $\omega=1$ and changes $L$ to
    $(\beta,1)\mathbin{\triangleright}L$, which is compatible with the
    trivial three-cochain by
    Lemma~\ref{thm:mpug_permuted_block_l_compatibility}. The first item of
    Lemma~\ref{thm:mpug_stabilizer_torsor} writes it as
    $(1,\gamma)\mathbin{\triangleright}L[\psi]$, so rescaling the action
    tensors by $\gamma^{-1}$ gives $L[\psi]$. For finite $G$, $\psi$ is
    cohomologous to its pointwise phase by
    Lemma~\ref{thm:mpug_circle_complex_units_h2}, and by
    Lemma~\ref{thm:mpug_stabilizer_cocycle_l_symbol} the coboundary is an
    action gauge, which is absorbed into the rescaling of the action tensors.
\end{proof}

\begin{lemma}[Block independence and extension of the stabilizer class]
    \label{lem:mpug_permuted_block_bi_extension}
    \lean{MPOTensor.GroupFamily.BlockActionData.isBlockIndependent_lSymbol_iff,
        MPOTensor.GroupFamily.BlockActionData.isBlockIndependent_lSymbol_iff_circle}
    \leanok
    \uses{lem:mpug_permuted_block_stabilizer_form,
        lem:mpug_permuted_block_l_gauge, lem:mpug_scalar_bi_restriction,
        lem:mpug_scalar_bi_restriction_circle}
    In the setting of Lemma~\ref{lem:mpug_permuted_block_stabilizer_form},
    take any rescaling of the tensors for which $L^x_{g_1,g_2}=\psi(h_1,h_2)$
    with a two-cocycle $\psi$ on $H$. The $L$-symbols are block independent
    after a joint scalar gauge if and only if $[\psi]\in H^2(H,\C^\times)$
    lies in the image of the restriction map from $H^2(G,\C^\times)$. If $G$
    is finite, this holds if and only if the class in $H^2(H,U(1))$ of the
    phase of $\psi$ lies in the image of the restriction map from
    $H^2(G,U(1))$.

    This is \cite[lines~7042--7064]{FrancoRubio2025MPUGauging} for the
    $L$-symbols of actual tensors, stated for the class of $\psi$; the
    representative form, an exact extension of $\psi$ to a cocycle on $G$, is
    Lemma~\ref{lem:mpug_scalar_bi_restriction}. The subgroup $H$ need not be
    normal.
\end{lemma}

\begin{proof}\leanok
    \uses{lem:mpug_permuted_block_l_gauge, lem:mpug_scalar_bi_restriction,
        lem:mpug_scalar_bi_restriction_circle}
    By Lemma~\ref{lem:mpug_permuted_block_l_gauge} the rescaled $L$-symbols
    are a joint scalar gauge of the original ones, and block independence is
    gauge invariant by Lemma~\ref{lem:mpug_scalar_bi_restriction_circle}.
    The two criteria for $L[\psi]$ are
    Lemmas~\ref{lem:mpug_scalar_bi_restriction}
    and~\ref{lem:mpug_scalar_bi_restriction_circle}.
\end{proof}

\subsection{Generalized cocycles of block-dependent cochains}

\begin{definition}[Generalized coboundary and cocycle]
    \label{def:mpug_generalized_cocycle}
    \lean{TNLean.Algebra.GeneralizedThreeCochain,
        TNLean.Algebra.GeneralizedFourCochain,
        TNLean.Algebra.ActionTensorGauge.coboundary,
        TNLean.Algebra.LSymbol.coboundary,
        TNLean.Algebra.GeneralizedThreeCochain.coboundary,
        TNLean.Algebra.LSymbol.IsGeneralizedCocycle}
    \leanok
    \uses{def:mpug_l_symbols, def:mpug_l_gauge}
    Let $\chi\colon G\times\mathsf X\to\C^\times$ be block-dependent. Its
    generalized coboundary is
    \begin{align}
        (d\chi)^x_{g,h}
         & =\frac{\chi^{h x}_g\chi^x_h}{\chi^x_{gh}}.
        \label{eq:mpug_generalized_coboundary}
    \end{align}
    For a block-dependent scalar $2$-cochain
    $f\colon\mathsf X\times G^2\to\C^\times$, set
    \begin{align}
        (df)^x_{g,h,k}
         & =\frac{f^x_{g,hk}f^x_{h,k}}
        {f^{k x}_{g,h}f^x_{gh,k}}.
        \label{eq:mpug_generalized_two_coboundary}
    \end{align}
    As recorded in \cite{gap:fbc25_two_cochain_coboundary_index_typo}, the
    final group index in the denominator is $k$, correcting the printed
    formula. For a block-dependent scalar $3$-cochain
    $\omega\colon\mathsf X\times G^3\to\C^\times$, the corresponding next
    coboundary is
    \begin{align}
        (d\omega)^x_{g,h,k,l}
         & =\frac{\omega^{l x}_{g,h,k}\omega^x_{g,hk,l}
            \omega^x_{h,k,l}}
        {\omega^x_{gh,k,l}\omega^x_{g,h,kl}}.
        \label{eq:mpug_generalized_three_coboundary}
    \end{align}
    A function $\Omega\colon\mathsf X\times G^2\to\C^\times$ is a
    generalized $2$-cocycle when
    \begin{align}
        \Omega^x_{g,hk}\Omega^x_{h,k}
         & =\Omega^{k x}_{g,h}\Omega^x_{gh,k}.
        \label{eq:mpug_generalized_cocycle}
    \end{align}
    The differential identities $d\Omega=1$ and $d(df)=1$ follow the group
    cohomology conventions of
    \cite[lines~5914--5922]{FrancoRubio2025MPUGauging}, with the
    three-to-four-cochain coboundary given explicitly by
    (\ref{eq:mpug_generalized_three_coboundary}).
    No transitivity or finiteness assumption is imposed on $\mathsf X$.

    % Local fix: docs/paper-gaps/fbc25_two_cochain_coboundary_index_typo.tex.
    % Source anchor: arXiv:2502.20257, lines 5914--5922 and lemma:G_cocycle.
\end{definition}

\begin{lemma}[Generalized coboundaries square to one]
    \label{thm:mpug_generalized_coboundary_squared}
    \lean{TNLean.Algebra.ActionTensorGauge.isGeneralizedCocycle_coboundary,
        TNLean.Algebra.GeneralizedThreeCochain.coboundary_coboundary}
    \leanok
    \uses{def:mpug_generalized_cocycle}
    The generalized coboundary of every block-dependent scalar $1$-cochain
    is a generalized $2$-cocycle. With the index correction of
    \cite{gap:fbc25_two_cochain_coboundary_index_typo}, every block-dependent
    scalar $2$-cochain $f$ satisfies the identity of
    \cite[line~5922]{FrancoRubio2025MPUGauging},
    \begin{align}
        d(df) & =1.
        \label{eq:mpug_generalized_coboundary_squared}
    \end{align}
\end{lemma}

\begin{proof}\leanok
    \uses{def:mpug_generalized_cocycle}
    Substituting the coboundary formulas, the action law
    $k(lx)=(kl)x$ and associativity identify the shifted block labels and
    group products. Every value of the original cochain then occurs once in
    the numerator and once in the denominator, so all factors cancel.
\end{proof}

\begin{definition}[Positive block-dependent cochains]
    \label{def:mpug_positive_generalized_cochains}
    \lean{TNLean.Algebra.PositiveUnits.embedding,
        TNLean.Algebra.PositiveUnits.IsPositive,
        TNLean.Algebra.ActionTensorGauge.IsPositiveValued,
        TNLean.Algebra.LSymbol.IsPositiveValued}
    \leanok
    \uses{def:mpug_l_symbols, def:mpug_l_gauge}
    Let $\iota\colon\mathbb R_{>0}\to\C^\times$ be the natural
    multiplicative embedding. A block-dependent scalar $1$-cochain $\chi$ is
    positive when
    \begin{align}
        \chi^x_g & \in\iota(\mathbb R_{>0})
        \label{eq:mpug_positive_one_cochain}
    \end{align}
    for every $x\in\mathsf X$ and $g\in G$. A block-dependent scalar
    $2$-cochain $\Omega$ is positive when
    \begin{align}
        \Omega^x_{g,h} & \in\iota(\mathbb R_{>0})
        \label{eq:mpug_positive_two_cochain}
    \end{align}
    for every $x\in\mathsf X$ and $g,h\in G$. Thus positivity means
    membership in the embedded subgroup of strictly positive real numbers,
    not positivity of the real part.
\end{definition}

\begin{definition}[Hat operation on block-dependent cochains]
    \label{def:mpug_generalized_hat}
    \lean{TNLean.Algebra.ActionTensorGauge.hat,
        TNLean.Algebra.LSymbol.hat}
    \leanok
    \uses{def:mpug_l_symbols, def:mpug_l_gauge}
    Following the one-index and two-index cases of the hat operation in
    \cite[line~5912]{FrancoRubio2025MPUGauging}, for block-dependent scalar
    cochains define
    \begin{align}
        \widehat\chi^{\,x}_g
         & =\chi^{\,g x}_{g^{-1}},
        \label{eq:mpug_generalized_hat_one}
    \end{align}
    and
    \begin{align}
        \widehat\Omega^{\,x}_{g,h}
         & =\Omega^{\,(gh)x}_{h^{-1},g^{-1}}.
        \label{eq:mpug_generalized_hat_two}
    \end{align}
\end{definition}

\begin{lemma}[Hat involution and generalized coboundaries]
    \label{thm:mpug_generalized_hat_coboundary}
    \lean{TNLean.Algebra.ActionTensorGauge.hat_hat,
        TNLean.Algebra.LSymbol.hat_hat,
        TNLean.Algebra.LSymbol.hat_coboundary}
    \leanok
    \uses{def:mpug_generalized_cocycle, def:mpug_generalized_hat}
    The block-dependent hat operation is an involution:
    \begin{align}
        \widehat{\widehat\chi} & =\chi,
        \label{eq:mpug_generalized_hat_hat_one}
    \end{align}
    and
    \begin{align}
        \widehat{\widehat\Omega} & =\Omega.
        \label{eq:mpug_generalized_hat_hat_two}
    \end{align}
    It commutes with the generalized coboundary:
    \begin{align}
        \widehat{d\chi} & =d\widehat\chi.
        \label{eq:mpug_generalized_hat_coboundary}
    \end{align}
\end{lemma}

\begin{proof}\leanok
    \uses{def:mpug_generalized_cocycle, def:mpug_generalized_hat}
    Applying (\ref{eq:mpug_generalized_hat_one}) twice and cancelling each
    group element with its inverse gives
    (\ref{eq:mpug_generalized_hat_hat_one}), and applying
    (\ref{eq:mpug_generalized_hat_two}) twice gives
    (\ref{eq:mpug_generalized_hat_hat_two}). For
    (\ref{eq:mpug_generalized_hat_coboundary}), the action law
    $g(h x)=(gh)x$ and $g^{-1}\bigl((gh)x\bigr)=h x$ give
    \begin{align}
        (d\widehat\chi)^x_{g,h}
         & =\frac{\widehat\chi^{\,h x}_g\widehat\chi^{\,x}_h}
        {\widehat\chi^{\,x}_{gh}}
        =\frac{\chi^{(gh)x}_{g^{-1}}\chi^{h x}_{h^{-1}}}
        {\chi^{(gh)x}_{(gh)^{-1}}}
        =(d\chi)^{(gh)x}_{h^{-1},g^{-1}}
        =\widehat{d\chi}^{\,x}_{g,h}.
    \end{align}
    The two sides therefore carry the same three values of $\chi$, reordered
    in $\C^\times$.
\end{proof}

\begin{definition}[Regular-action matrices]
    \label{def:mpug_generalized_regular_matrices}
    \lean{TNLean.Algebra.LSymbol.regularMatrix}
    \leanok
    \uses{def:mpug_l_symbols}
    Let $\Omega\colon\mathsf X\times G^2\to\C^\times$ be an $L$-symbol.
    As in \cite[Lemma~\texttt{lemma:G\_cocycle}]{FrancoRubio2025MPUGauging},
    define the regular-action matrix $X^x_g$ on the group algebra $\C G$ by
    \begin{align}
        X^x_g\ket{k}
         & =\Omega^{k^{-1}x}_{g,k}\ket{gk}.
        \label{eq:mpug_generalized_regular_matrix}
    \end{align}
    % Source anchor: arXiv:2502.20257, lines 5938--5942.
\end{definition}

\begin{lemma}[Generalized projective multiplication law]
    \label{thm:mpug_generalized_regular_multiplication}
    \lean{TNLean.Algebra.LSymbol.regularMatrix_mul}
    \leanok
    \uses{def:mpug_generalized_cocycle,
        def:mpug_generalized_regular_matrices}
    Let $G$ be finite. If $\Omega$
    satisfies~(\ref{eq:mpug_generalized_cocycle}), then
    \begin{align}
        X^{h x}_gX^x_h
         & =\Omega^x_{g,h}X^x_{gh}.
        \label{eq:mpug_generalized_regular_multiplication}
    \end{align}
    % Source anchor: arXiv:2502.20257, lemma:G_cocycle, lines 5931--5943,
    %   which states this multiplication law "for G being a finite group".
    %   The blueprint statement previously omitted that hypothesis, which the
    %   formal statement carries.
\end{lemma}

\begin{proof}\leanok
    \uses{def:mpug_generalized_cocycle,
        def:mpug_generalized_regular_matrices}
    Applying the left side to $\ket{k}$ yields the coefficient
    \begin{align}
        \Omega^{(hk)^{-1}h x}_{g,hk}\Omega^{k^{-1}x}_{h,k}
         & =\Omega^{k^{-1}x}_{g,hk}\Omega^{k^{-1}x}_{h,k} \\
         & =\Omega^x_{g,h}\Omega^{k^{-1}x}_{gh,k},
    \end{align}
    where the last equality is
    (\ref{eq:mpug_generalized_cocycle}) at the block $k^{-1}x$.
    Formula~(\ref{eq:mpug_generalized_regular_matrix}) now gives
    (\ref{eq:mpug_generalized_regular_multiplication}).
\end{proof}

\begin{lemma}[Finite-group power trivialization]
    \label{thm:mpug_generalized_cocycle_power}
    \lean{TNLean.Algebra.LSymbol.regularMatrix_eq_permutation_mul_diagonal,
        TNLean.Algebra.LSymbol.regularDeterminant_ne_zero,
        TNLean.Algebra.LSymbol.determinantCochain,
        TNLean.Algebra.LSymbol.pow_card_eq_coboundary_apply,
        TNLean.Algebra.LSymbol.pow_card_eq_coboundary}
    \leanok
    \uses{def:mpug_generalized_cocycle,
        def:mpug_generalized_regular_matrices}
    Let $G$ be finite and let $\Omega$ be a generalized $2$-cocycle. Then
    $\Omega^{|G|}$ is a generalized coboundary. More precisely, with
    \begin{align}
        \chi^x_g
         & =\det X^x_g,
        \label{eq:mpug_generalized_determinant_cochain}
    \end{align}
    one has
    \begin{align}
        \left(\Omega^x_{g,h}\right)^{|G|}
         & =\frac{\det X^{h x}_g\det X^x_h}{\det X^x_{gh}}
        =(d\chi)^x_{g,h}.
        \label{eq:mpug_generalized_power_coboundary}
    \end{align}
    This establishes
    \cite[Lemma~\texttt{lemma:G\_cocycle}]{FrancoRubio2025MPUGauging}; as
    recorded in \cite{gap:fbc25_generalized_cocycle_determinant_typo}, the
    second numerator determinant in
    (\ref{eq:mpug_generalized_power_coboundary}) carries the index $h$,
    correcting the printed formula.

    % Local fix: docs/paper-gaps/fbc25_generalized_cocycle_determinant_typo.tex.
    % Source anchor: arXiv:2502.20257, lemma:G_cocycle, lines 5931--5948.
\end{lemma}

\begin{proof}\leanok
    \uses{def:mpug_generalized_regular_matrices,
        thm:mpug_generalized_regular_multiplication}
    Let $P_g$ be the regular-action permutation matrix defined by
    $P_g\ket{k}=\ket{gk}$, and set
    $D^x_g=\operatorname{diag}(\Omega^{k^{-1}x}_{g,k})_{k\in G}$.
    Then
    \begin{align}
        X^x_g
         & =P_gD^x_g.
    \end{align}
    Since $\det P_g\in\{1,-1\}$ and every diagonal entry of $D^x_g$
    is nonzero, $\det X^x_g\ne0$. Thus
    (\ref{eq:mpug_generalized_determinant_cochain}) takes values in
    $\C^\times$. Taking determinants in
    (\ref{eq:mpug_generalized_regular_multiplication}), using
    multiplicativity of the determinant on the left and
    $\det(cX)=c^{|G|}\det X$ for the $|G|\times|G|$ matrix $X^x_{gh}$ on the
    right, gives
    \begin{align}
        \det X^{h x}_g\det X^x_h
         & =\left(\Omega^x_{g,h}\right)^{|G|}\det X^x_{gh}.
        \label{eq:mpug_generalized_determinant_product}
    \end{align}
    Dividing (\ref{eq:mpug_generalized_determinant_product}) by
    $\det X^x_{gh}$ gives
    (\ref{eq:mpug_generalized_power_coboundary}).
\end{proof}

\begin{lemma}[Positive generalized cocycles are coboundaries]
    \label{thm:mpug_positive_generalized_cocycle}
    \lean{TNLean.Algebra.PositiveUnits.positiveRpow,
        TNLean.Algebra.PositiveUnits.positiveRpow_inv_natCast_pow,
        TNLean.Algebra.ActionTensorGauge.positiveRpow,
        TNLean.Algebra.LSymbol.positivePrimitive,
        TNLean.Algebra.LSymbol.isPositiveValued_positivePrimitive,
        TNLean.Algebra.LSymbol.exists_positive_coboundary}
    \leanok
    \uses{def:mpug_generalized_cocycle,
        def:mpug_positive_generalized_cochains,
        thm:mpug_generalized_cocycle_power}
    Let $G$ be finite and let $G$ act on a set $\mathsf X$. If $\Omega$ is a
    positive generalized $2$-cocycle, then there is a positive
    block-dependent scalar $1$-cochain $\chi_0$ such that
    \begin{align}
        \Omega^x_{g,h}
         & =(d\chi_0)^x_{g,h}
        =\frac{\chi^{\,h x}_{0,g}\chi^x_{0,h}}
        {\chi^x_{0,gh}}
        \label{eq:mpug_positive_coboundary}
    \end{align}
    for all $x\in\mathsf X$ and $g,h\in G$.
    This establishes the first assertion of
    \cite[Corollary at lines~5950--5960]{FrancoRubio2025MPUGauging}; as
    recorded in \cite{gap:fbc25_positive_generalized_cocycle_block_index},
    whereas the source writes $\chi_g$,
    Lemma~\ref{thm:mpug_generalized_cocycle_power} constructs $\chi^x_g$
    through (\ref{eq:mpug_generalized_determinant_cochain}) and the block label
    $x$ is therefore retained.
    % Local correction:
    % docs/paper-gaps/fbc25_positive_generalized_cocycle_block_index.tex.
    % Source anchor: arXiv:2502.20257, cor:positive_trivial, lines 5950--5956.
\end{lemma}

\begin{proof}\leanok
    \uses{def:mpug_positive_generalized_cochains,
        thm:mpug_generalized_cocycle_power}
    With $\delta^x_g=\det X^x_g$ denoting the determinant cochain from
    Lemma~\ref{thm:mpug_generalized_cocycle_power}, set
    \begin{align}
        \chi^x_{0,g}
         & =\lvert\delta^x_g\rvert^{1/\lvert G\rvert}.
        \label{eq:mpug_positive_determinant_root}
    \end{align}
    As recorded in
    \cite{gap:fbc25_positive_generalized_cocycle_determinant_sign}, because
    the regular-action permutation can contribute a sign to $\delta^x_g$, its
    norm is taken before the positive $\lvert G\rvert$-th root.
    Since $\Omega$ is positive and the complex norm is multiplicative, taking
    norms in (\ref{eq:mpug_generalized_power_coboundary}) gives
    \begin{align}
        \Omega^{\lvert G\rvert} & =d\lvert\delta\rvert.
        \label{eq:mpug_positive_norm_power}
    \end{align}
    Positive real powers are multiplicative, so
    (\ref{eq:mpug_positive_determinant_root}) and
    (\ref{eq:mpug_positive_norm_power}) give
    \begin{align}
        (d\chi_0)^{\lvert G\rvert}
         & =d\lvert\delta\rvert=\Omega^{\lvert G\rvert},
        \label{eq:mpug_positive_root_power}
    \end{align}
    and the $\lvert G\rvert$-th power is injective on $\mathbb R_{>0}$. Hence
    $d\chi_0=\Omega$.
    % Consolidated clarification:
    % docs/paper-gaps/fbc25_positive_generalized_cocycle_determinant_sign.tex.
\end{proof}

\begin{lemma}[Hat-compatible positive primitive]
    \label{thm:mpug_positive_generalized_cocycle_hat}
    \lean{TNLean.Algebra.ActionTensorGauge.coboundary_div,
        TNLean.Algebra.ActionTensorGauge.coboundary_positiveRpow,
        TNLean.Algebra.LSymbol.exists_positive_coboundary_hat_mul_eq_one}
    \leanok
    \uses{def:mpug_generalized_hat,
        def:mpug_positive_generalized_cochains,
        thm:mpug_generalized_hat_coboundary,
        thm:mpug_positive_generalized_cocycle}
    Under the hypotheses of
    Lemma~\ref{thm:mpug_positive_generalized_cocycle}, suppose additionally
    that
    \begin{align}
        \Omega^x_{g,h}\widehat\Omega^{\,x}_{g,h}
         & =1
        \label{eq:mpug_positive_hat_hypothesis}
    \end{align}
    for every $x\in\mathsf X$ and $g,h\in G$. Then, as in the second assertion
    of \cite[Corollary at lines~5950--5960]{FrancoRubio2025MPUGauging}, the
    positive cochain can be chosen so that
    \begin{align}
        \Omega & =d\chi,
        \label{eq:mpug_positive_hat_coboundary}
    \end{align}
    and
    \begin{align}
        \chi^x_g\widehat\chi^{\,x}_g
         & =1.
        \label{eq:mpug_positive_hat_witness_identity}
    \end{align}
    % Source anchor: arXiv:2502.20257, cor:positive_trivial, lines 5955--5960.
\end{lemma}

\begin{proof}\leanok
    \uses{thm:mpug_generalized_hat_coboundary,
        thm:mpug_positive_generalized_cocycle}
    Let $\chi_0$ satisfy $d\chi_0=\Omega$ as in
    Lemma~\ref{thm:mpug_positive_generalized_cocycle}. Define
    \begin{align}
        \chi^x_g
         & =\sqrt{\frac{\chi^x_{0,g}}
            {\widehat\chi^{\,x}_{0,g}}}
        =\sqrt{\frac{\chi^x_{0,g}}
            {\chi^{\,g x}_{0,g^{-1}}}}.
        \label{eq:mpug_positive_hat_witness}
    \end{align}
    All values are positive, so the positive square roots are unambiguous.
    Using (\ref{eq:mpug_generalized_hat_coboundary}) and the hypothesis
    (\ref{eq:mpug_positive_hat_hypothesis}) gives
    \begin{align}
        d\chi
         & =\sqrt{d\chi_0/d\widehat\chi_0}
        =\sqrt{\Omega/\widehat\Omega}
        =\sqrt{\Omega^2}
        =\Omega.
    \end{align}
    The same substitution in
    (\ref{eq:mpug_positive_hat_witness}) gives
    \begin{align}
        \widehat\chi
         & =\sqrt{\widehat\chi_0/\chi_0}
        =\chi^{-1},
    \end{align}
    which proves
    (\ref{eq:mpug_positive_hat_witness_identity}).
\end{proof}

\subsection{Inversion cochains and the group of order two}

\begin{definition}[Inversion cochains]
    \label{def:mpug_inversion_cochains}
    \lean{TNLean.Algebra.ScalarOneCochain,
        TNLean.Algebra.ScalarOneCochain.coboundary,
        TNLean.Algebra.ScalarCocycle.hat,
        TNLean.Algebra.ScalarThreeCochain.hat,
        TNLean.Algebra.ScalarThreeCochain.Xi,
        TNLean.Algebra.ScalarThreeCochain.sigma}
    \leanok
    \uses{def:mpug_scalar_cochains}
    These conventions and inversion cochains follow
    \cite[lines~5910--5916 and 5967--5973]{FrancoRubio2025MPUGauging}.
    For scalar cochains with one, two, and three arguments, set
    \begin{align}
        (d\sigma)(g,h) & =\frac{\sigma(g)\sigma(h)}{\sigma(gh)}.
        \label{eq:mpug_one_coboundary}
    \end{align}
    The hat reverses the arguments and replaces each group element by its
    inverse:
    \begin{align}
        \widehat\beta(g,h) & =\beta(h^{-1},g^{-1}).
        \label{eq:mpug_two_hat}
    \end{align}
    \begin{align}
        \widehat\omega(g,h,k) & =\omega(k^{-1},h^{-1},g^{-1}).
        \label{eq:mpug_three_hat}
    \end{align}
    Define the inversion cochains by
    \begin{align}
        \Xi(g,h) & =
        \frac{\omega(h^{-1},g^{-1},g)}
        {\omega(h^{-1}g^{-1},g,h)}.
        \label{eq:mpug_intro_Xi}
    \end{align}
    \begin{align}
        \sigma(g) & =\omega(g,g^{-1},g).
        \label{eq:mpug_intro_sigma}
    \end{align}
    % Source anchors: arXiv:2502.20257, lines 5910--5916 and eq:introXi.
\end{definition}

\begin{lemma}[The hat operation is an involution]
    \label{thm:mpug_scalar_hat_involution}
    \lean{TNLean.Algebra.ScalarCocycle.hat_hat,
        TNLean.Algebra.ScalarThreeCochain.hat_hat}
    \leanok
    \uses{def:mpug_inversion_cochains}
    Let $G$ be an arbitrary group. Every scalar $2$-cochain $\beta$ and every
    scalar $3$-cochain $\omega$ satisfy
    \begin{align}
        \widehat{\widehat\beta} & =\beta,
        \label{eq:mpug_two_hat_hat}
    \end{align}
    \begin{align}
        \widehat{\widehat\omega} & =\omega.
        \label{eq:mpug_three_hat_hat}
    \end{align}
    % Source anchor: arXiv:2502.20257, line 5912, where the hat operation is
    %   introduced as an involution for cochains of any arity.
\end{lemma}

\begin{proof}\leanok
    \uses{def:mpug_inversion_cochains}
    Reversing the arguments of (\ref{eq:mpug_two_hat}) and inverting each of
    them twice restores the original arguments, and likewise for
    (\ref{eq:mpug_three_hat}).
\end{proof}

\begin{lemma}[Two inversion-cochain formulas]
    \label{thm:mpug_inversion_cochain_formulas}
    \lean{TNLean.Algebra.ScalarThreeCochain.Xi_eq_second_formula,
        TNLean.Algebra.ScalarThreeCochain.sigma_inv}
    \leanok
    \uses{def:mpug_scalar_cochains, def:mpug_three_cocycle,
        def:mpug_inversion_cochains}
    These formulas are part of the cocycle calculation in
    \cite[lines~5967--5979]{FrancoRubio2025MPUGauging}.
    Let $G$ be an arbitrary group and let $\omega$ be a normalized scalar
    $3$-cocycle. Then
    \begin{align}
        \Xi(g,h) & =
        \frac{\omega(h^{-1},g^{-1},g)}
        {\omega(h^{-1}g^{-1},g,h)}
        =
        \frac{\omega(h^{-1},g^{-1},gh)}
        {\omega(g^{-1},g,h)}.
        \label{eq:mpug_Xi_two_formulas}
    \end{align}
    The inversion one-cochain satisfies
    \begin{align}
        \sigma(g^{-1}) & =\sigma(g)^{-1}.
        \label{eq:mpug_sigma_inverse}
    \end{align}
    % Source anchor: arXiv:2502.20257, eq:introXi and lemma:omegaXi.
\end{lemma}

\begin{proof}\leanok
    \uses{def:mpug_scalar_cochains, def:mpug_three_cocycle,
        def:mpug_inversion_cochains}
    Applying (\ref{eq:mpug_three_cocycle}) to $(h^{-1},g^{-1},g,h)$ and using
    (\ref{eq:mpug_three_normalized}) gives the second quotient in
    (\ref{eq:mpug_Xi_two_formulas}). Applying
    (\ref{eq:mpug_three_cocycle}) to
    $(g,g^{-1},g,g^{-1})$ gives
    $1=\sigma(g)\sigma(g^{-1})$, hence
    (\ref{eq:mpug_sigma_inverse}).
\end{proof}

\begin{lemma}[Inversion identities for a normalized three-cocycle]
    \label{thm:mpug_omega_Xi}
    \lean{TNLean.Algebra.ScalarThreeCochain.hat_eq_mul_coboundary_Xi,
        TNLean.Algebra.ScalarThreeCochain.Xi_eq_hat_mul_coboundary_sigma}
    \leanok
    \uses{def:mpug_scalar_cochains, def:mpug_three_cocycle,
        def:mpug_fusion_gauge, def:mpug_inversion_cochains}
    These inversion identities and their cocycle substitutions are given in
    \cite[lines~5967--5979]{FrancoRubio2025MPUGauging}.
    Let $\omega$ be a normalized scalar $3$-cocycle. Then
    \begin{align}
        \widehat\omega & =\omega\,d\Xi.
        \label{eq:mpug_hat_omega}
    \end{align}
    The inversion two-cochain also satisfies
    \begin{align}
        \Xi & =\widehat\Xi\,d\sigma.
        \label{eq:mpug_Xi_hat_sigma}
    \end{align}
    % Source anchor: arXiv:2502.20257, lemma:omegaXi, lines 5967--5979.
\end{lemma}

\begin{proof}\leanok
    \uses{def:mpug_three_cocycle, def:mpug_fusion_gauge,
        def:mpug_inversion_cochains, thm:mpug_inversion_cochain_formulas}
    For (\ref{eq:mpug_hat_omega}), multiplying the cocycle identities at
    \begin{align}
         & (k^{-1},h^{-1},g^{-1},g),\quad
        (k^{-1},h^{-1}g^{-1},g,h),\quad
        (k^{-1}h^{-1}g^{-1},g,h,k),
    \end{align}
    and cancelling the normalized factors gives the equality. For
    (\ref{eq:mpug_Xi_hat_sigma}), substituting the second expression in
    (\ref{eq:mpug_Xi_two_formulas}) for $\widehat\Xi$ and combining the
    identities at
    \begin{align}
         & (h^{-1}g^{-1},g,h,h^{-1}g^{-1}),\quad
        (h^{-1},h,h^{-1},g^{-1}),\quad
        (h^{-1},g^{-1},g,g^{-1}),\quad
        (g,g^{-1},g,g^{-1})
    \end{align}
    gives (\ref{eq:mpug_sigma_inverse}) in the three required places at $g$,
    $h$, and $gh$.
\end{proof}

Let $\Z_2=\{e,s\}$ with $s^2=e$.

\begin{lemma}[Normalized $L$-symbol ratio]
    \label{thm:mpug_sigma_l_ratio}
    \lean{TNLean.Algebra.LSymbol.sigma_eq_div}
    \leanok
    \uses{def:mpug_l_symbols, def:mpug_l_normalized,
        def:mpug_inversion_cochains}
    Let $G$ act on a set $\mathsf X$. Suppose that a normalized $L$-symbol
    $L$ is compatible with a scalar $3$-cochain $\omega$. Then, as in
    \cite[lines~5327--5329]{FrancoRubio2025MPUGauging}, for every
    $x\in\mathsf X$ and $g\in G$,
    \begin{align}
        \sigma(g)
         & =\frac{L^x_{g^{-1},g}}{L^{g x}_{g,g^{-1}}}.
        \label{eq:mpug_sigma_l_ratio}
    \end{align}
    % Source anchor: arXiv:2502.20257, lines 5327--5329.
\end{lemma}

\begin{proof}\leanok
    \uses{def:mpug_l_symbols, def:mpug_l_normalized,
        def:mpug_inversion_cochains}
    Applying the compatibility relation
    (\ref{eq:mpug_l_compatibility}) to $(g,g^{-1},g)$, the two factors with an
    identity argument are $1$ by normalization, leaving
    (\ref{eq:mpug_sigma_l_ratio}).
\end{proof}

\begin{lemma}[Trivial scalar gauge class for $\Z_2$]
    \label{thm:mpug_z2_trivial_scalar_class}
    \lean{TNLean.Algebra.ScalarThreeCochain.isTrivialGaugeClass_iff_sigma_generator_eq_one}
    \leanok
    \uses{def:mpug_scalar_cochains, def:mpug_fusion_gauge,
        def:mpug_trivial_gauge_class, def:mpug_inversion_cochains}
    Let $\omega$ be a normalized scalar $3$-cochain on $\Z_2$. Its gauge
    class is trivial if and only if
    \begin{align}
        \sigma(s) & =1.
        \label{eq:mpug_z2_trivial_sigma}
    \end{align}
    This strengthens the scalar step in
    \cite[lines~5339--5341]{FrancoRubio2025MPUGauging}: only normalization,
    rather than the cocycle identity, is needed for this equivalence.
\end{lemma}

\begin{proof}\leanok
    \uses{def:mpug_fusion_gauge,
        thm:mpug_gauge_preserves_normalization,
        thm:mpug_cohomologous_to_equivalence}
    Suppose first that $\omega=d\beta$ for an arbitrary scalar $2$-cochain
    $\beta$. Although $\beta$ need not be normalized, the normalization of
    $d\beta$ implies that there is a constant $c\in\C^\times$ such that
    $\beta(e,h)=\beta(h,e)=c$ for every $h\in\Z_2$. Since $s^2=e$,
    \begin{align}
        (d\beta)(s,s,s)
         & =\frac{\beta(s,e)\beta(s,s)}
        {\beta(s,s)\beta(e,s)}=\frac{c}{c}=1.
        \label{eq:mpug_z2_coboundary_generator}
    \end{align}
    Thus every scalar $2$-cochain whose coboundary is normalized contributes
    the factor $1$ at $(s,s,s)$. In particular, an arbitrary gauge between
    normalized scalar $3$-cochains cannot change this value, and
    $\sigma(s)=1$. Conversely, normalization determines every
    value of $\omega$ except $\omega(s,s,s)=\sigma(s)$. Hence
    (\ref{eq:mpug_z2_trivial_sigma}) makes $\omega$ the constant cochain $1$.
\end{proof}

\begin{lemma}[Sign dichotomy for the $\Z_2$ inversion scalar]
    \label{thm:mpug_z2_sigma_sign}
    \lean{TNLean.Algebra.ScalarThreeCochain.sigma_generator_eq_one_or_eq_neg_one}
    \leanok
    \uses{def:mpug_three_cocycle, def:mpug_inversion_cochains}
    Let $\omega$ be a normalized scalar $3$-cocycle on $\Z_2$. Then
    \begin{align}
        \sigma(s) & \in\{1,-1\}.
        \label{eq:mpug_z2_sigma_sign}
    \end{align}
\end{lemma}

\begin{proof}\leanok
    \uses{thm:mpug_inversion_cochain_formulas}
    Since $s^{-1}=s$, formula (\ref{eq:mpug_sigma_inverse}) gives
    $\sigma(s)=\sigma(s)^{-1}$. Hence $\sigma(s)^2=1$, whose two solutions in
    $\C^\times$ are $1$ and $-1$.
\end{proof}

For an element $g$ whose square is the identity, the source condition
$\xi_g(x)\xi_{g^{-1}}(g\mathbin{\cdot}x)=-1$ reduces to
$\xi_g(x)\xi_g(g\mathbin{\cdot}x)=-1$ because $g^{-1}=g$.

\begin{definition}[Anomalous sign for an involution]
    \label{def:mpug_anomalous_involution_sign}
    \lean{TNLean.Algebra.IsComplexUnitSign,
        TNLean.Algebra.HasAnomalousInvolutionSign}
    \leanok
    Let $f\colon\mathsf X\to\mathsf X$. An anomalous sign for $f$ is a
    function $\xi\colon\mathsf X\to\C^\times$ such that
    \begin{align}
        \xi(x) & \in\{1,-1\},
        \label{eq:mpug_anomalous_sign_range}
    \end{align}
    and
    \begin{align}
        \xi(x)\xi(f(x)) & =-1
        \label{eq:mpug_anomalous_sign_pairing}
    \end{align}
    for every $x\in\mathsf X$.
    % Source anchor: arXiv:2502.20257, prop:def_tens_anomalous,
    % eq:defects_Z2b, lines 2979--2981.
\end{definition}

\begin{lemma}[Anomalous signs and fixed points]
    \label{thm:mpug_anomalous_involution_sign_iff}
    \lean{TNLean.Algebra.hasAnomalousInvolutionSign_of_fixedPointFree,
        TNLean.Algebra.not_hasAnomalousInvolutionSign_of_fixedPoint,
        TNLean.Algebra.hasAnomalousInvolutionSign_iff_fixedPointFree,
        TNLean.Algebra.involutive_smul_of_mul_self_eq_one,
        TNLean.Algebra.hasAnomalousInvolutionSign_smul_iff}
    \leanok
    \uses{def:mpug_anomalous_involution_sign}
    Let $f\colon\mathsf X\to\mathsf X$ be an involution. An anomalous sign
    exists if and only if $f$ has no fixed point. In particular, if a group
    acts on $\mathsf X$ from the left and $g^2=e$, the sign condition for
    anomalous defect tensors in
    \cite[lines~2979--2981]{FrancoRubio2025MPUGauging} holds for the action
    $x\mapsto g\mathbin{\cdot}x$ if and only if
    $g\mathbin{\cdot}x\neq x$ for every $x\in\mathsf X$.
    % Source anchor: arXiv:2502.20257, prop:def_tens_anomalous,
    % eq:defects_Z2b, lines 2979--2981.
\end{lemma}

\begin{proof}\leanok
    \uses{def:mpug_anomalous_involution_sign}
    Choosing a well-order on $\mathsf X$ and assigning $1$ to the smaller
    element and $-1$ to the larger in each two-element orbit $\{x,f(x)\}$,
    involutivity exchanges the two elements so that their signs multiply to
    $-1$. Conversely, at a fixed point the product is the square of either
    $1$ or $-1$, hence equals $1$ rather than $-1$. For the action statement,
    $g\mathbin{\cdot}(g\mathbin{\cdot}x)=(g^2)\mathbin{\cdot}x=x$.
\end{proof}

\begin{lemma}[Nontrivial scalar gauge class for $\Z_2$]
    \label{thm:mpug_z2_nontrivial_scalar_class}
    \lean{TNLean.Algebra.ScalarThreeCochain.not_isTrivialGaugeClass_iff_sigma_generator_eq_neg_one}
    \leanok
    \uses{def:mpug_three_cocycle, def:mpug_trivial_gauge_class,
        def:mpug_inversion_cochains}
    For a normalized scalar $3$-cocycle $\omega$ on $\Z_2$, the scalar
    cohomology classification used in the proof at
    \cite[lines~5339--5341]{FrancoRubio2025MPUGauging} states that its gauge
    class is nontrivial if and only if
    \begin{align}
        \sigma(s) & =-1.
        \label{eq:mpug_z2_nontrivial_sigma}
    \end{align}
    % Source anchor: arXiv:2502.20257, lines 5339--5341.
\end{lemma}

\begin{proof}\leanok
    \uses{thm:mpug_z2_trivial_scalar_class, thm:mpug_z2_sigma_sign}
    By Lemma~\ref{thm:mpug_z2_trivial_scalar_class}, a nontrivial class has
    $\sigma(s)\neq1$. The dichotomy
    (\ref{eq:mpug_z2_sigma_sign}) then forces $\sigma(s)=-1$, and the converse
    follows from the same two statements.
\end{proof}

This completes the classification of the normalized scalar $3$-cocycle by
$\sigma(s)$ in the setting of \cite[lines~5207--5338]{FrancoRubio2025MPUGauging}.

